\documentclass[11pt]{article}
\usepackage[margin=1.1in]{geometry}
\usepackage{amsmath,amssymb,amsthm}
\usepackage{booktabs}
\usepackage{lmodern}
\usepackage{microtype}
\usepackage[colorlinks=true,linkcolor=blue,citecolor=blue,urlcolor=blue]{hyperref}

\newtheorem{conjecture}{Conjecture}
\newtheorem{remark}{Remark}
\theoremstyle{definition}
\newtheorem{question}{Question}

\newcommand{\Li}{\operatorname{Li}}
\newcommand{\e}{\mathrm{e}}
\newcommand{\SG}{\mathfrak S}
\newcommand{\SGc}{\mathfrak S^{\mathrm{c}}}
\newcommand{\PD}{\mathrm{PD}}
\newcommand{\Exp}{\mathrm{Exp}}
\DeclareMathOperator{\Var}{Var}
\DeclareMathOperator{\Cov}{Cov}
\DeclareMathOperator{\Hom}{Hom}
\DeclareMathOperator{\Lie}{Lie}
\DeclareMathOperator{\Gal}{Gal}
\DeclareMathOperator{\Aut}{Aut}
\DeclareMathOperator{\fix}{fix}
\DeclareMathOperator{\rank}{rank}
\DeclareMathOperator{\codim}{codim}
\DeclareMathOperator{\disc}{disc}

\title{Twenty-five structural conjectures for prime patterns, occupancy,
finite logarithms, arboreal dynamics and factorization}
\date{August 2026}

\begin{document}
\maketitle

\begin{abstract}
Twenty-five conjectures are stated across five programmes: connected
prime-pattern fields, arithmetic occupancy and local-global decoupling,
finite logarithms and Kummer rank, arboreal resolution and primitive
dynamical factors, and Frobenius-marked polynomial factorization.  Each
statement declares its probability space, separates its constructive
layer from its converse layer, derives its constants from local, Galois,
overlap or sieve structure or else registers the derivation as a stated
programme, and carries a finite-range evidence note that reports what the
computation can and cannot bear on.  The organizing principle is that a
conjecture should introduce a mechanism, a canonical invariant, a new
scale or a functorial relationship, so that a generic limit law by itself
is not admitted.
\end{abstract}

\section*{Summary sheet: the twenty-five statements}
\addcontentsline{toc}{section}{Summary sheet}

\begin{enumerate}\itemsep2pt

\item \textbf{Connected Hardy--Littlewood generating functional
(Conjecture~\ref{c1}).}  Every joint cumulant of centred motif fields,
under the window-position sampling law, expands over the overlap poset
with M\"obius-connected singular-series coefficients, no coefficient
occurring outside that expansion, and Gaussianity follows from a stated
bound on connected cumulants of order at least three rather than being
assumed.

\item \textbf{Overlap-filtration variance law (Conjecture~\ref{c2}).}
The variance of a linear combination of motif counts expands over the
two-level overlap poset of a pair, with explicit powers of $\log X$
indexed by coincidence multiplicity, annihilation of the coincidence
matrices deletes that whole level, and the renormalized disjoint kernel
then fixes the leading scale.

\item \textbf{Local-to-spectral covariance decomposition
(Conjecture~\ref{c3}).}  The motif covariance splits into a coincidence
part indexed by the overlap poset, a renormalized connected
singular-series kernel of disjoint translates, and the window pullback of
the joint zero-correlation measure, with no fourth source at leading
order and a crossover scale fixed by equality of masses.

\item \textbf{Spectrally conditioned extremal index
(Conjecture~\ref{c4}).}  On the window-position probability space,
conditioning on the sigma-field generated by the long-wavelength
component makes cluster starts Poisson, the within-cluster law is the
admissible self-overlap graph, and removing the conditioning returns a
Cox cluster process.

\item \textbf{Completeness of arithmetic wavelets
(Conjecture~\ref{c5}).}  Joint kernels of the overlap incidence tensors
produce predicted low-variance combinations, and the converse layer
asserts that every anomalously rigid coefficient sequence is
asymptotically one of them.

\item \textbf{First-order arithmetic occupancy kernel
(Conjecture~\ref{c6}).}  Under the modulus-class law the least-prime
survival function is $\e^{-t}[1-A(t)/\log Q+o(1/\log Q)]$ with $A$
obtained by integrating the connected two-prime exclusion kernel.

\item \textbf{Dual additive-multiplicative spectral geometry
(Conjecture~\ref{c7}).}  The least-prime covariance operator is an
additive convolution kernel plus a character-diagonal operator built from
zeros of Dirichlet $L$-functions, with no surviving mixed term after
dyadic averaging.

\item \textbf{Conditioned Cox--Poisson law for the last classes
(Conjecture~\ref{c8}).}  On the modulus-class probability space,
conditioning on the finitely many non-vanishing low modes Poissonizes the
terminal excess process, and unconditionally the limit is the Cox process
directed by those modes.

\item \textbf{Rank-one exceptional-zero deformation
(Conjecture~\ref{c9}).}  A real zero close to $1$ deforms the first-arrival
field by exactly one spectral spike along the associated real character.

\item \textbf{Growing local-condition independence for odd class groups
(Conjecture~\ref{c10}).}  After the two-primary genus data are fixed, the
odd class group remains Cohen--Lenstra distributed under prescribed
splitting signs at all primes up to $B$ whenever $\pi(B)=o(\log|D|)$.

\item \textbf{Kummer--Haar law for Fermat quotients
(Conjecture~\ref{c11}).}  For a saturated group of rank $r$ the normalized
Fermat-quotient homomorphism equidistributes on the relation subtorus,
whose dimension is the saturated rank and not the number of displayed
generators.

\item \textbf{Square-root large sieve for finite logarithms
(Conjecture~\ref{c12}).}  Nonrelation character sums of finite logarithms
over dyadic prime blocks are of square-root size uniformly for character
heights growing as a power of the range.

\item \textbf{Saturated-rank threshold for simultaneous Wieferich primes
(Conjecture~\ref{c13}).}  Rank one gives $\kappa(\Gamma)\log\log X$
vanishing primes and rank at least two gives finitely many, the
invariance statements holding only for saturated and Kummer-compatible
groups.

\item \textbf{Rank-depth Poisson law for higher strata
(Conjecture~\ref{c14}).}  Over a uniformly sampled height ensemble of
generator tuples, the maximal-depth exceptional primes form a Poisson
process of intensity $\kappa(r,k;p)p^{-rk}-\kappa(r,k+1;p)p^{-r(k+1)}$,
each fixed depth slice is Poisson of intensity $\kappa(r,k;p)p^{-rk}$,
and the depths are joined by Haar lifts.

\item \textbf{Toric finite logarithm and codimension principle
(Conjecture~\ref{c15}).}  The first congruence quotient on any algebraic
torus is Haar on the Lie algebra of the algebraic subgroup generated by
the point, functorially in the torus, so the exceptional exponent is
$\dim S$.

\item \textbf{Joint arboreal Chebotarev law at growing depth
(Conjecture~\ref{c16}).}  Below the entropy-conductor boundary the whole
visible preimage tree is Haar--Chebotarev distributed in the finite
arboreal image, jointly across levels.

\item \textbf{Sharp entropy-conductor barrier (Conjecture~\ref{c17}).}
The upper half is a conditional proposition under a Riemann hypothesis for
the preimage fields, the lower half is asserted family-averaged, and the
critical-window Gaussian-multinomial profile is the conjectural clause.

\item \textbf{Poisson--Dirichlet law for primitive dynamical factors
(Conjecture~\ref{c18}).}  In a uniformly sampled parameter family with
slowly growing level, the ranked logarithms of the primitive prime
factors converge to $\PD(1)$, marked by first-appearance level and
Frobenius data.

\item \textbf{Euler law for squarefull primitive factors
(Conjecture~\ref{c19}).}  The primitive part is squarefree with a limiting
Euler-product probability whose local factors count parameters modulo
$p^2$, and the squarefull logarithm is tight.

\item \textbf{Geometric-arboreal criterion for large dynamical gcds
(Conjecture~\ref{c20}).}  Macroscopic orbit gcds occur if and only if the
orbits lie on a periodic algebraic correspondence, the ``only if''
direction being the converse layer.

\item \textbf{Frobenius-marked Poisson--Dirichlet factor process
(Conjecture~\ref{c21}).}  Macroscopic prime factors of polynomial values
form a $\PD(1)$ process marked by Frobenius classes with intensity
size-biased by fixed-point counts, normalized to one by Burnside.

\item \textbf{Small-large independence after logarithmic mass removal
(Conjecture~\ref{c22}).}  The small-prime Kubilius field is asymptotically
independent of the large-factor process once the latter is normalized by
the residual mass rather than the full mass.

\item \textbf{Coloured factor superposition with local-exclusion gluing
(Conjecture~\ref{c23}).}  Conditioned on the joint local state the colour
processes are independent $\PD(1)$, and unconditionally their entire
interaction is an explicit Gibbs--Euler factor built from joint root-count
tensors.

\item \textbf{Saddle-point entanglement law for joint smoothness
(Conjecture~\ref{c24}).}  The joint smooth-value density is the product of
component probabilities times an Euler product of local root
distributions tilted by the multivariate saddle-point parameters, so the
entanglement factor moves with the smoothness parameters.

\item \textbf{Buchstab--Bateman--Horn sieve-to-prime bridge
(Conjecture~\ref{c25}).}  Rough-value densities follow the finite-$u$
Buchstab law, and the conditional prime probability is the Bateman--Horn
density divided by that law, with an elementary endpoint at $1<u\le2$.

\end{enumerate}

\tableofcontents

\section{Notation and conventions}\label{sec:notation}

\subsection*{Sampling laws}

Every stochastic clause below names one of three ensembles.  No clause
attaches ``almost surely'', a central limit theorem, a conditional law or
an extremal limit to a deterministic sequence.

\emph{(W) Window-position law.}  Fix $X\ge3$, a fixed $\varepsilon>0$ and
a window length $L$ with $X^{\varepsilon}\le L\le X^{1-\varepsilon}$.  Put
$\Omega_X=[X,2X]$ with its Borel field and the dyadic-log probability
\[
  \mathbb P_X(dx)\;=\;\frac{dx}{x\log2}.
\]
A \emph{random window} is $[x,x+L]$ with $x\sim\mathbb P_X$.  All fields
of Programme~I are random variables on $(\Omega_X,\mathbb P_X)$, and every
sigma-field conditioned on in Conjectures~\ref{c4} and~\ref{c8} is a
sub-sigma-field of a space declared here.  Expectation, variance and
covariance under $\mathbb P_X$ are written $\mathbb E_X$, $\Var_X$,
$\Cov_X$.

\emph{(M) Modulus-class law.}  Fix $Q\ge3$ and put
\[
  \Omega_Q=\{(q,a):\ q\ \text{prime},\ Q\le q\le2Q,\ a\in(\mathbb
  Z/q\mathbb Z)^{\times}\},\qquad
  \mathbb P_Q(q,a)=\frac1{\pi(2Q)-\pi(Q)}\cdot\frac1{q-1}.
\]
The modulus is uniform on the dyadic block of primes and the class is
uniform given the modulus.  This is the ensemble meant by ``dyadic
averaging'' throughout Programme~II.

\emph{(P) Parameter-family laws.}  In Programme~III, Conjecture~\ref{c14}
samples an ordered $r$-tuple of multiplicatively independent integers of
height at most $H$ uniformly from the admissible tuples, admissibility
excluding the Kummer-incompatible tuples defined below.  In
Programme~IV, Conjectures~\ref{c18}, \ref{c19} sample the parameter $c$
uniformly from the integers with $0<|c|\le H$ lying outside the
postcritically finite and reducible loci, and Conjecture~\ref{c17} uses
the same law for its family-averaged clause.  In Programme~V the integer
$n$ is uniform on $\{1,\dots,N\}$.  Conjecture~\ref{c16} samples $p$
uniformly from the primes $p\le X$ that are unramified in the relevant
field.

\subsection*{Prime motifs, the overlap poset and connected singular series}

Throughout, $p$ and $q$ denote primes, $\pi(x)$ is the prime counting
function, $\gamma=0.5772156\ldots$ is Euler's constant,
$\Li(x)=\int_2^{x}dt/\log t$, and $\Lambda$ is the von Mangoldt function.
For an algebraic point, $h$ denotes the absolute logarithmic Weil height
and $\hat h_f$ the canonical height of a polynomial $f$.  We write
$\mathbb E_X$, $\Var_X$, $\Cov_X$ for moments under the window-position
law and $\mathbb E_Q$ for expectation under the modulus-class law.

A \emph{motif} is a finite set $H=\{h_1<\dots<h_k\}\subset\mathbb
Z_{\ge0}$ with $h_1=0$.  It is \emph{admissible} if for every prime $p$
the reduction of $H$ modulo $p$ omits a residue class.  With
$\nu_H(p)=\#(H\bmod p)$ the singular series is
\[
  \SG(H)=\prod_p\frac{1-\nu_H(p)/p}{(1-1/p)^{k}} .
\]
Write $\Lambda_H(n)=\prod_{h\in H}\Lambda(n+h)$.  For a smooth compactly
supported $w$ and $x\in\Omega_X$ the \emph{centred weighted motif field}
is
\[
  Y(H;x,L)\;=\;\sum_{n\in\mathbb Z}w\!\left(\frac{n-x}{L}\right)
  \bigl(\Lambda_H(n)-\SG(H)\bigr).
\]
Under (W) this is a random variable, and $V=V(H;X,L)=\Var_XY(H;\cdot,L)$
is its \emph{second connected mass}.

Let $H_1,\dots,H_m$ be motifs.  A \emph{configuration} of order $r$ is a
choice of indices $i_1,\dots,i_r$ together with shifts
$t_1,\dots,t_r\in\mathbb Z$.  Its \emph{overlap type} is the set of
coincidences
\[
  \tau=\bigl\{\{(u,h),(v,h')\}:\ u\neq v,\ h\in H_{i_u},\ h'\in H_{i_v},
  \ h+t_u=h'+t_v\bigr\},
\]
recorded up to simultaneous translation.  Overlap types are ordered by
inclusion and form the \emph{overlap poset}
$\mathcal O(H_{i_1},\dots,H_{i_r})$, whose minimum is the empty type (all
translates disjoint) and whose grading $\codim\tau$ is the number of
independent linear relations among $t_1,\dots,t_r$ that $\tau$ imposes.
The empty type is called the \emph{disjoint stratum}.

For a finite family $A$ of translated motifs let $\SG(A)$ be the singular
series of the union pattern $\bigcup_{a\in A}a$, defined whenever that
union is admissible.  The \emph{connected singular series} is the
M\"obius inversion of $\SG$ over the partition lattice $\Pi(A)$:
\[
  \SGc(A)\;=\;\sum_{\pi\in\Pi(A)}\mu(\pi,\hat1)\prod_{B\in\pi}\SG(B),
  \qquad \mu(\pi,\hat1)=(-1)^{|\pi|-1}(|\pi|-1)!,
\]
so that $\SGc(A)$ vanishes whenever $\SG$ factorizes over a proper
partition of $A$.  This is the exact sense of ``connected'' used
throughout Programme~I, and it is what isolates genuine $r$-way
arithmetic interaction from products of lower-order interactions.

For $r\ge2$ and $j\ge0$ the \emph{incidence tensor} $A^{(j)}_r$ of the
family is the symmetric $r$-th order form on $\mathbb R^{m}$ whose entry
at $(i_1,\dots,i_r)$ is the total connected singular-series mass of the
overlap types of codimension $j$ carried by the $r$-fold configurations
of $H_{i_1},\dots,H_{i_r}$, weighted by the window integrals of
Conjecture~\ref{c1}.  The order of the tensor is the cumulant order $r$,
whereas the superscript $j$ indexes the strata by codimension and is not
a tensor order, so the two indices never coincide in meaning.  In
particular at second order $A^{(j)}:=A^{(j)}_2$ is a symmetric $m\times
m$ matrix and $A^{(j)}c$ is a vector, for every $j$ and not only for
$j=2$.  Two translates carry
a single shift parameter, so the overlap poset of a pair has exactly two
levels, codimension $0$, the disjoint stratum, in which the shift is
free, and codimension $1$, the coincidence strata, in which a coincidence
pins the shift to one of finitely many values.  Thus $j\in\{0,1\}$ when
$r=2$.  The coincidence strata are graded further by the coincidence
multiplicity $e(\tau)\ge1$ of Conjecture~\ref{c1}, and $A^{(1,e)}$
denotes the part of $A^{(1)}$ carried by the types with $e(\tau)=e$, so
that the coincidence matrix of scale $X$ is
\[
  A^{(1)}(X)\;=\;\sum_{e\ge1}(\log X)^{e}A^{(1,e)} .
\]
The disjoint stratum $j=0$ carries no constant tensor.  Its mass is the
renormalized kernel matrix $K_{\mathrm{HL}}(L)$ of
Conjecture~\ref{c3}, which is a symmetric $m\times m$ matrix for each $L$
but depends on the actual shift values and on $L$ rather than on the
overlap type alone, and
it is placed last in the filtrations of Conjectures~\ref{c2}
and~\ref{c5} because the annihilation conditions are imposed on the
coincidence tensors first.  ``Last stratum'' therefore refers to the
order in which the strata are examined and not to the codimension
grading, in which the disjoint stratum is the minimum.

\subsection*{Least primes and occupancy}

For prime $q$ and $a\in(\mathbb Z/q\mathbb Z)^{\times}$ let $p(q,a)$ be
the least prime congruent to $a$ modulo $q$ and set
$T_q(a)=\Li(p(q,a))/(q-1)$.  Under (M) this is a random variable.  The
\emph{first-arrival transform} $\mathcal F$ has for domain the bounded
measurable kernels $K\colon\mathbb R_{\ge0}\to\mathbb R$ and for codomain
the continuous functions on $\mathbb R_{\ge0}$ vanishing to second order
at the origin, and it is defined by the linearization
\[
  \mathcal F[K](t)\;=\;\tfrac12\int_0^{t}\!\!\int_0^{t}K(|u-v|)\,du\,dv
  \;=\;\int_0^{t}(t-u)K(u)\,du .
\]
This is the derivative at the constant unit intensity of the map that
sends the connected two-point correlation of a point process on
$\Li$-time to minus the logarithm of its first-arrival survival function:
a process of unit intensity and connected pair correlation $K$ has
$-\log\mathbb P[\text{no arrival in }[0,t]]=t-\mathcal
F[K](t)+O(\|K\|_\infty^{2})$.  The \emph{connected two-prime exclusion
kernel} is the family of numbers $\SG(\{0,kq\})-1$ for $k\ge1$, where
$\SG(\{0,d\})$ is the singular series of the pair motif of shift $d$,
these being the connected densities of two primes lying in the same class
modulo $q$.  Extending $T_q$ to $a=0$ by the recorded value at the zero
class, the \emph{additive spectrum} and \emph{multiplicative spectrum} of
the centred field $\widetilde T_q(a)=T_q(a)-\mathbb E_a[T_q]$ are
\[
  S^{\mathrm{add}}_q(\xi)=\frac1q\Bigl|\sum_{a\bmod q}\widetilde T_q(a)
  \exp(2\pi ia\xi/q)\Bigr|^2,
  \qquad
  S^{\mathrm{mult}}_q(\chi)=\frac1{q-1}\Bigl|\sum_{a\neq0}\widetilde
  T_q(a)\overline{\chi(a)}\Bigr|^2 ,
\]
for $\xi\in\mathbb Z/q\mathbb Z$ and $\chi$ a Dirichlet character modulo
$q$.  Both are deterministic for each $q$ and become random variables
under (M).

\subsection*{Class groups}

For a negative fundamental discriminant $D$ write $h(D)$ for the class
number of $\mathbb Q(\sqrt D)$ and $\mathrm{Cl}(D)$ for the class group.
The \emph{two-primary genus data} of $D$ is the pair consisting of the
genus character group of $D$, equivalently the set of prime discriminant
factors of $D$, together with the induced two-primary quotient of
$\mathrm{Cl}(D)$.  For an odd prime $\ell$, the \emph{Cohen--Lenstra law}
for $\mathrm{Cl}(D)[\ell^{\infty}]$ assigns a finite abelian $\ell$-group
$G$ the mass proportional to $1/|\Aut G|$.

\subsection*{Finite logarithms, saturation and Kummer compatibility}

For a prime $p$ and $a\in\mathbb Z$ with $p\nmid a$ the \emph{Fermat
quotient} is $q_p(a)=\bigl((a^{p-1}-1)/p\bigr)\bmod p$.  For a finitely
generated $\Gamma\le\mathbb Q^{\times}$ with $p$ dividing no numerator or
denominator of a generator, $q_p$ extends to a homomorphism
$q_p\colon\Gamma\to\mathbb F_p$.

$\Gamma$ is \emph{saturated} if $\Gamma=\{x\in\mathbb Q^{\times}:x^{n}\in
\Gamma\ \text{for some}\ n\ge1\}$, and $\Gamma^{\mathrm{sat}}$ denotes
the saturation.  Saturation does not change $\Gamma\otimes\mathbb Q$, so
$\rank\Gamma=\rank\Gamma^{\mathrm{sat}}=:r$.

The \emph{Kummer degree} of $\Gamma$ at level $m$ is
\[
  \Delta_\Gamma(m)=[\mathbb Q(\zeta_m,\Gamma^{1/m}):\mathbb Q(\zeta_m)],
\]
a divisor of $m^{r}$.  $\Gamma$ is \emph{Kummer-generic at $p$} if
$\Delta_\Gamma(p)=p^{r}$, and the \emph{local Kummer factor} is
\[
  \kappa(\Gamma;p)\;=\;\frac{p^{r}}{\Delta_\Gamma(p)}\;\ge\;1,
\]
equal to $1$ exactly at the Kummer-generic primes.  A tuple of generators
is \emph{Kummer-compatible} if the group it generates is saturated and
Kummer-generic at all but finitely many primes.  Since $(\mathbb
Z/p^2\mathbb Z)^{\times}$ is cyclic of order $p(p-1)$, its subgroup of
$p$-th powers is exactly $\{x:x^{p-1}=1\}$, so
\[
  q_p(a)=0\iff a^{p-1}\equiv1\ (\mathrm{mod}\ p^2)
  \iff a\bmod p^2\ \text{is a $p$-th power},
\]
which is the identification of $q_p$ with the connecting map $\Gamma\to
(\mathbb Z/p^2\mathbb Z)^{\times}/\bigl((\mathbb Z/p^2\mathbb
Z)^{\times}\bigr)^{p}\cong\mathbb Z/p\mathbb Z$ used throughout
Programme~III.

Characters are taken on the ambient generator torus and not on $\Gamma$.
Fix an ordered generating tuple $g=(g_1,\dots,g_s)$ of $\Gamma$, so that
the Fermat-quotient vector
\[
  \mathbf q_p(g)/p\;=\;\bigl(q_p(g_1),\dots,q_p(g_s)\bigr)/p
  \ \in\ (\mathbb R/\mathbb Z)^{s}
\]
lives on the ambient torus $(\mathbb R/\mathbb Z)^{s}$.  Let
\[
  R(g)\;=\;\Bigl\{m\in\mathbb Z^{s}:\ \prod_{i=1}^{s}g_i^{m_i}=1\Bigr\}
\]
be the \emph{relation lattice} of the tuple, of rank $s-r$, let
$R(g)^{\mathrm{sat}}=(R(g)\otimes\mathbb Q)\cap\mathbb Z^{s}$ be its
saturation in $\mathbb Z^{s}$, and let $\mathbb T_\Gamma\le(\mathbb
R/\mathbb Z)^{s}$ be the identity component of the annihilator of
$R(g)$, a subtorus of dimension exactly $r$.  A \emph{character} is an
element $m\in\mathbb Z^{s}$, acting on the ambient torus by $\langle
m,\mathbf q_p(g)\rangle=\sum_im_iq_p(g_i)$, and its \emph{height} is
$\max_i|m_i|$.  A character $m$ is a \emph{nonrelation character} of the
tuple if $m\notin R(g)^{\mathrm{sat}}$, equivalently if $m$ restricts
nontrivially to $\mathbb T_\Gamma$.  The exclusion is exactly the right
one, since for $m\in R(g)^{\mathrm{sat}}$, say with $km\in R(g)$, one has
$k\langle m,\mathbf q_p(g)\rangle\equiv0\pmod p$ and hence $\langle
m,\mathbf q_p(g)\rangle\equiv0\pmod p$ at every $p\nmid k$, so such a
character is constant on the values and carries no information, whereas
for $m\in\Hom(\Gamma,\mathbb Z)$ the relations of $\Gamma$ are killed
automatically and no nontrivial exclusion is available at all.

\subsection*{Arboreal notation}

Let $f\in\mathbb Q[x]$ have degree $d\ge2$ and let $a\in\mathbb Q$ have
infinite backward orbit.  Write $K_n=\mathbb Q(f^{-n}(a))$ for the $n$-th
preimage field, $G_n=\Gal(K_n/\mathbb Q)$, $D_n=\disc(K_n)$, and $T_n$
for the tree of preimages of $a$ to level $n$, so that
$G_n\hookrightarrow\Aut(T_n)=[S_d]^{\wr n}$, the $n$-fold iterated wreath
product of the symmetric group $S_d$, of order
$(d!)^{(d^{n}-1)/(d-1)}$.  Only for $d=2$ is this an iterated wreath
product of cyclic groups.  The \emph{arboreal image}
is $G_\infty=\varprojlim G_n\le\Aut(T_\infty)$, and it has \emph{finite
index} if $[\Aut(T_\infty):G_\infty]<\infty$.  A \emph{conjugacy
cylinder} at level $n$ is a conjugacy class of $G_n$, and a
\emph{cylinder event} is the corresponding set of factorization types of
$f^{j}(x)-a$ modulo $p$ for $j\le n$.  The \emph{entropy-conductor
functional} is
\[
  E(n,X)\;=\;\frac{|G_n|}{\pi(X)}
  \;+\;\frac{|G_n|^2\bigl(\log D_n+[K_n:\mathbb Q]\log X\bigr)^2}{X},
\]
the first term being the information-theoretic requirement that every
state of $G_n$ receive enough primes and the second the square-root
Chebotarev conductor obstruction.  A polynomial is \emph{disintegrated}
if it is not conjugate to a power map, a Chebyshev polynomial or a
Latt\`es map, this being the class in which the Medvedev--Scanlon
classification of invariant subvarieties applies.

For $f_c(x)=x^{d}+c$ put $v_n(c)=f_c^{n}(0)$.  The \emph{primitive part}
$P_n(c)$ is the largest divisor of $|v_n(c)|$ composed only of primes
dividing no $v_m(c)$ with $1\le m<n$.

\subsection*{Factor processes and Poisson--Dirichlet marks}

$\PD(1)$ denotes the Poisson--Dirichlet law of parameter $1$, the law of
the decreasing rearrangement of the normalized jumps of a gamma
subordinator, equivalently the limit law of the ranked ratios $\log
p_i/\log n$ over the prime factors of a uniform integer $n\le N$.  Its
largest coordinate has mean the Golomb--Dickman constant
$\lambda=0.6243299\ldots$.  Given a mark space $S$ and a probability
kernel $\nu(\cdot\mid u)$ on $S$, a \emph{marked $\PD(1)$ process} is the
point process $\sum_i\delta_{(u_i,s_i)}$ on $(0,1]\times S$ in which
$(u_i)$ is $\PD(1)$ and the marks $s_i$ are conditionally independent
given $(u_i)$ with $s_i\sim\nu(\cdot\mid u_i)$.  For irreducible $f$ with
splitting group $G$ acting on the roots, $\fix(C)$ is the number of roots
fixed by any element of the conjugacy class $C$, and Burnside's lemma
gives $\sum_C(|C|/|G|)\fix(C)=1$ for transitive $G$, so the mark
intensity of Conjecture~\ref{c21} needs no fitted normalization.

For a polynomial value $f(n)$ and a parameter $y$, the \emph{$y$-smooth
part} $S_y(f(n))$ is the largest $y$-smooth divisor and
$R_y=f(n)/S_y(f(n))$ is the \emph{residual cofactor}.  The
\emph{small-prime Kubilius field} is the sigma-field generated by the
valuations $v_p(f(n))$ for $p\le y$.  For irreducible $f$ we write
$\rho_f(p)=\#\{x\bmod p:p\mid f(x)\}$, $V_f(y)=\prod_{p\le
y}(1-\rho_f(p)/p)$, $C_f=\prod_p(1-\rho_f(p)/p)(1-1/p)^{-1}$ for the
Bateman--Horn constant, and $\omega$ for the Buchstab function, the
continuous solution of $(u\omega(u))'=\omega(u-1)$ with $\omega(u)=1/u$
on $1<u\le2$.  For several polynomials, $\rho(I,p)=\#\{x\bmod p:p\mid
g_i(x)\ \text{for all}\ i\in I\}$ is the \emph{joint root-count tensor}.

\subsection*{Tagging conventions}

Clauses marked \textbf{[elementary]} are proved by the argument indicated
with them.  Clauses marked \textbf{[conditional proposition]} are
deductions from a hypothesis named in place.  Clauses marked
\textbf{[strong converse layer]} are exhaustiveness or completeness
assertions.  These are stated separately from the constructive clauses of
the same conjecture and are separately falsifiable: a counterexample to a
converse clause leaves the constructive clause of that conjecture
standing, and the two therefore carry independent risk.  Clauses marked
\textbf{[stated programme]} record the derivation route for a constant or
kernel that is canonical but not yet evaluated in closed form.  All
remaining clauses are conjectural.  Questions raised in the
\textsc{Question} environment are open problems and are not counted among
the twenty-five conjectures.

\subsection*{Evidence conventions}

The evidence note attached to each conjecture reports a finite-range
diagnostic and nothing more.  Computations are consistency checks, never
support for an asymptotic claim.  Quoted standard errors and $z$-scores
are descriptive, and where samples are correlated, overlapping or
selected after inspection they are not calibrated tail probabilities.
Where a range is too short to identify an asymptotic parameter this is
said.  All prime-pattern window statistics below were sampled at step
$L/4$ within a window of length $L$, so consecutive samples overlap in
three quarters of their support and the effective number of independent
samples is about one quarter of the nominal count, with no error bars
computed by the generating script.  The generating script samples window
positions on a uniform grid over $[X,2X]$ rather than from the
dyadic-log density of (W), an admissible discretization of (W) because
the two densities differ by a factor bounded between $1/(2\log2)$ and
$1/\log2$ on that interval, so no reported statistic is affected beyond
that bounded reweighting.  Reproducible bounds are primes to
$3{,}000{,}000$, prime moduli below $1000$, Fermat quotients to
$250{,}000$, imaginary quadratic discriminants to $15{,}000$, and the
polynomial and dynamical samples recorded with each note.

\section{Programme I. Connected prime-pattern fields}\label{sec:I}

The programme builds a connected generating functional for prime motifs,
an overlap filtration of its variance, a decomposition of its covariance
into local and spectral halves, a conditioned extreme-value theory, and a
converse classification of low-variance combinations.  All five
statements live on the window-position probability space (W) of
Section~\ref{sec:notation}.

\begin{conjecture}[Connected Hardy--Littlewood generating
functional]\label{c1}
Fix admissible motifs $H_1,\dots,H_m$, a smooth compactly supported
window $w$, and the window-position law (W).  Write
$Y_i=Y(H_i;\cdot,L)$ and let $\kappa_r$ denote the joint cumulant of
order $r$ under $\mathbb P_X$.  Then:
\emph{(i)} [overlap expansion] for every $r\ge2$ and every choice of
indices $i_1,\dots,i_r$ the cumulant
$\kappa_r\bigl(Y_{i_1},\dots,Y_{i_r}\bigr)$ has an asymptotic expansion
indexed by the overlap poset $\mathcal O(H_{i_1},\dots,H_{i_r})$, in
which each stratum carries an explicit power of $\log X$ and a
renormalized kernel sum over its free shifts.  For $\tau\in\mathcal O$
put $\mathrm{fr}(\tau)=r-1-\codim\tau$ for the number of free shifts of
$\tau$, and for $s\in\mathbb Z^{\mathrm{fr}(\tau)}$ let $U(\tau;s)$ be
the union pattern of the configuration of type $\tau$ with free shifts
$s$, a set of distinct integers, and let $\SGc(\tau;s)=\SGc\bigl(U(\tau;
s)\bigr)$ be its connected singular series, obtained by M\"obius
inversion over the partition lattice as in Section~\ref{sec:notation}.
Let
\[
  e(\tau)\;=\;\sum_{u=1}^{r}|H_{i_u}|\;-\;|U(\tau;s)|
\]
be the \emph{coincidence multiplicity} of $\tau$, a quantity independent
of $s$, so that a point of the union pattern carried by $\mu$ translates
contributes $\mu-1$ to $e(\tau)$.  Then
\[
  \kappa_r\bigl(Y_{i_1},\dots,Y_{i_r}\bigr)
  \;=\;\sum_{\tau\in\mathcal O}(\log X)^{e(\tau)}\,\mathcal
  K_\tau(L)\;+\;\text{lower strata},
\]
\[
  \mathcal K_\tau(L)\;=\;\sum_{s\in\mathbb Z^{\mathrm{fr}(\tau)}}
  w_L^{(\tau)}(s)\,\SGc(\tau;s),\qquad
  w_L^{(\tau)}(s)=L\int_{\mathbb R}w(t)\prod_{j=1}^{\mathrm{fr}(\tau)}
  w\!\left(t+\frac{s_j}{L}\right)dt ,
\]
the window weight $w_L^{(\tau)}$ contributing one factor $L$ for the
summation over $n$ and one for each free shift, so that $\mathcal
K_\tau(L)$ sits at scale $L^{\,r-\codim\tau}$ up to the logarithmic
renormalization described in clause (v).  The power of $\log X$ is the
coincidence bookkeeping: a point of multiplicity $\mu$ contributes
$\Lambda^{\mu}$, and $\sum_{n\le x}\Lambda(n)^{\mu}\sim x(\log
x)^{\mu-1}$.  In the smallest case $H=\{0\}$ and $r=2$ the
full-coincidence stratum has $e(\tau)=1$ and $\codim\tau=1$ and
contributes $L\log X$, matching $\sum_{n\le x}\Lambda(n)^{2}\sim x\log
x$.
\emph{(ii)} [strong converse layer] no coefficient and no scale occurs
outside this expansion: for every $s\ge0$ the difference between
$\kappa_r$ and the expansion of (i) truncated at codimension $s$ is
$o\bigl(L^{\,r-s}(\log X)^{e_{\max}}\bigr)$, where $e_{\max}$ is the
largest coincidence multiplicity occurring at codimension at most $s$.
Clause (ii) is separately falsifiable from (i), since a persistent
normalized cumulant not assigned to any overlap stratum refutes (ii)
while leaving the expansion of (i) intact.
\emph{(iii)} [hypothesis $\mathrm H_r$] let $V=V(H_i;X,L)$ and let
$\kappa_r^{\mathrm{disj}}=\mathcal K_{\hat0}(L)$ denote the contribution
of the disjoint stratum, the renormalized kernel sum of (i) at
$\codim\tau=0$.  The hypothesis $\mathrm H_r$, for $r\ge3$, asserts that
$\kappa_r^{\mathrm{disj}}=o\bigl(V^{r/2}\bigr)$ in every range in which
$V\to\infty$.  This is a hypothesis and not a consequence of (i) and
(ii), which fix the shape of the expansion but not the size of its
disjoint stratum.
\emph{(iv)} [conditional proposition, under $\mathrm H_r$ for every
$r\ge3$] $Y_i/\sqrt V$ converges in law under $\mathbb P_X$ to the
standard Gaussian, and all higher joint cumulants of the normalized
fields obey Wick pairing.  Gaussianity is thus the conclusion drawn from
clause (iii) and is not an independent input to the conjecture.
\emph{(v)} [renormalization of the free-shift strata] the free-shift
strata, and in particular the disjoint stratum $\codim\tau=0$, are
asserted only in the kernel-sum form displayed in (i), and not as a
constant depending on $\tau$ alone times a window integral.  The union
pattern, and therefore $\SGc(\tau;s)$, depends on the shift values $s$
and not merely on the type, and the centred shift sum diverges
logarithmically, the two-motif case being
\[
  \sum_{0<h\le H}\bigl(\SG(\{0,h\})-1\bigr)\;=\;-\tfrac12\log H+O(1),
\]
in the form established by Montgomery and Soundararajan.  Consequently
$\mathcal K_\tau(L)$ is smaller than $L^{\,r-\codim\tau}$ by a
logarithmic factor on the free-shift strata, and the kernel of the
disjoint stratum at $r=2$ is exactly the renormalized kernel
$K_{\mathrm{HL}}$ of Conjecture~\ref{c3}.
\end{conjecture}

M\"obius inversion over set partitions is what isolates genuine $r$-way
arithmetic interaction, since $\SGc$ vanishes on any family whose
singular series factorizes, while the overlap poset records the equal
prime constraints that a partition cannot see.  Together they play the
role of connected diagrams in a cluster expansion, and the statement
turns Hardy--Littlewood constants into the coefficients of a stochastic
field theory for prime constellations rather than a list of separate
predictions.  If proved it would compute joint moments of motif counts,
identify the exact obstruction to Gaussianity through clause (iii), and
supply Conjectures~\ref{c2} to~\ref{c5} with a common parent instead of a
family of analogies.  The nearest antecedents are Gallagher's averaging
of singular series, the moment method of Montgomery and Soundararajan,
the singular-series estimates of Pintz, and the odd-moment work of
Kuperberg, each of which treats a projection of the object above rather
than the full multitype functional.

\emph{Evidence.}  For four prime-pair fields at primes below
$3\times10^{6}$ the standardized fourth connected cumulants lay between
$-0.31$ and $0.41$ and the selected mixed fourth cumulants between
$-0.13$ and $0.26$ as the window mean rose from $2.5$ to $93$, but
windows were sampled at step $L/4$ and therefore overlap in three
quarters of their support, so about one quarter of the nominal samples
are effective, no error bars accompany any of these numbers, and the
range identifies no coefficient of the expansion.

\begin{conjecture}[Overlap-filtration variance law]\label{c2}
Fix admissible motifs $H_1,\dots,H_m$, a window $w$ and the law (W).  For
$c\in\mathbb R^{m}$ put $Y_c=\sum_ic_iY(H_i;\cdot,L)$.  Then:
\emph{(i)} [variance filtration] $\Var_XY_c$ has an asymptotic expansion
indexed by the overlap poset of the pairs $(H_i,H_j)$, which at second
order has exactly the two levels recorded in
Section~\ref{sec:notation}.  Grading the codimension-one strata by the
coincidence multiplicity $e(\tau)\ge1$ of Conjecture~\ref{c1},
\[
  \Var_XY_c\;=\;\sum_{e\ge1}L(\log X)^{e}\,c^{\mathsf T}A^{(1,e)}c
  \;+\;c^{\mathsf T}K_{\mathrm{HL}}(L)\,c\;+\;o(L),
\]
in which $A^{(1,e)}$ is the symmetric $m\times m$ incidence matrix of
Section~\ref{sec:notation} carrying the connected singular-series masses
of the codimension-one overlap types of coincidence multiplicity $e$
together with their window integrals, and $K_{\mathrm{HL}}(L)$ is the
renormalized disjoint kernel matrix of Conjecture~\ref{c3},
$K_{\mathrm{HL}}(i,j)=\sum_hw_L(h)\SGc(H_i,H_j+h)$.  The disjoint stratum
enters as a kernel sum over the actual shifts and not as a constant times
$L^{2}$, since the centred shift sum converges only after the
logarithmic renormalization of Conjecture~\ref{c1}(v), and it carries no
divergent mass.
\emph{(ii)} [annihilation, of depth one at second order] if
$A^{(1,e)}c=0$ for every $e\ge1$ then the whole coincidence level
vanishes from the expansion and the leading scale of $\Var_XY_c$ is
$c^{\mathsf T}K_{\mathrm{HL}}(L)c$, provided that quantity does not
vanish.  Because $j\in\{0,1\}$ at $r=2$, the annihilation chain has depth
one: a second-order statistic admits exactly one nontrivial annihilation
step, and any deeper cancellation is a condition on the higher-order
tensors $A^{(j)}_r$ with $r\ge3$ and is not visible in the variance
alone.
\emph{(iii)} [strong converse layer] no variance scale occurs outside the
coincidence matrices $A^{(1,e)}$ and the renormalized disjoint kernel,
separately falsifiable from (i) and (ii) in the sense that an unassigned
scale refutes (iii) alone.
\end{conjecture}

Two occurrences of a motif may share $0,1,2,\dots$ prime constraints, and
each sharing pattern contributes at its own power of the window length
and its own power of $\log X$, so the overlap lattice is the arithmetic
analogue of a diagram expansion and the variance is graded by it.  The consequence is constructive:
balanced observables can be designed to annihilate the low-codimension
diagrams and thereby manufacture rigid prime statistics, which is the
route to the arithmetic wavelets classified in Conjecture~\ref{c5}, and
it would explain why apparently similar prime races fluctuate at
different scales.  Tuple correlations in the tradition of Hardy and
Littlewood already distinguish coincident from disjoint shifts, and
Gallagher's and Kuperberg's averaging arguments handle the extremes of
the filtration, but a general variance grading indexed by the overlap
poset is the content claimed here.

\emph{Evidence.}  The covariance matrix of four pair-pattern fields kept
its three smaller eigenvalues at $0.29$ to $0.44$ of the largest with
condition number near $3.4$ across window lengths $400$ to $15{,}000$,
which is consistent with a persistent mode hierarchy rather than a single
common trend but identifies no stratum, and the overlapping-window
caveat again applies, the step $L/4$ leaving about a quarter as many
effective samples and no error bars.

\begin{conjecture}[Local-to-spectral covariance decomposition for prime
motifs]\label{c3}
Fix admissible motifs $H_i$, a smooth compactly supported window $w$, a
fixed $\varepsilon>0$ and the law (W), and let $X^{\varepsilon}\le L\le
X^{1-\varepsilon}$.  Fix a modulus $Q_0=Q_0(H_i,H_j)$ divisible by every
prime $p\le|H_i|+|H_j|$ and restrict $n$ to a fixed class modulo $Q_0$,
so that each translated von Mangoldt factor $\Lambda(n+h)$ is expanded by
the explicit formula in an arithmetic progression modulo $Q_0$.  Let
$\mathcal L(H)$ be the resulting finite set of Dirichlet $L$-functions
$L(s,\chi)$, $\chi$ modulo $Q_0$, one for each character occurring in
that expansion for the motif $H$, and write $\mathcal Z(H)$ for the
multiset of nontrivial zeros of the members of $\mathcal L(H)$, each zero
carried with the coefficient $\overline{\chi(a+h)}/\varphi(Q_0)$ with
which it enters.  Let
\[
  M_{w,L,x}(s)\;=\;\int_0^{\infty}w\!\left(\frac{u-x}{L}\right)u^{\,s-1}\,du
\]
be the window Mellin transform, so that the explicit formula gives
\[
  Y(H;x,L)\;=\;-\sum_{\rho\in\mathcal Z(H)}M_{w,L,x}(\rho)
  \;+\;\text{local terms} .
\]
Let $\mu_{ij}$ be the \emph{joint zero-correlation measure} of the pair,
that is, the pair-correlation measure of the two zero multisets,
\[
  \mu_{ij}\;=\;\sum_{\rho\in\mathcal Z(H_i)}\ \sum_{\rho'\in\mathcal
  Z(H_j)}c_\rho\overline{c_{\rho'}}\,
  \delta_{(\rho,\rho')} ,
\]
$c_\rho$ denoting the coefficient carried by $\rho$.  Then:
\emph{(i)} [decomposition] the motif covariance admits a unique
asymptotic decomposition
\[
  \Cov_X\bigl(Y(H_i;\cdot,L),Y(H_j;\cdot,L)\bigr)\;=\;
  D_{\mathrm{ov}}(i,j)+K_{\mathrm{HL}}(i,j)+K_{\mathrm{spec}}(i,j)+o(V),
\]
in which $D_{\mathrm{ov}}$ is the sum over the strata of positive
codimension of the overlap poset of the pair, carrying the same
logarithmic bookkeeping as Conjecture~\ref{c1},
\[
  D_{\mathrm{ov}}(i,j)\;=\;\sum_{e\ge1}L(\log X)^{e}\,A^{(1,e)}_{ij},
\]
$K_{\mathrm{HL}}$ is the renormalized connected kernel of the disjoint
stratum, given by
\[
  K_{\mathrm{HL}}(i,j)=\sum_{h}w_L(h)\,\SGc(H_i,H_j+h),\qquad
  w_L(h)=L\int_{\mathbb R}w(t)w(t+h/L)\,dt ,
\]
with $w_L$ the window autocorrelation weight, and $K_{\mathrm{spec}}$ is
the pullback of $\mu_{ij}$ through the window Mellin transforms,
\[
  K_{\mathrm{spec}}(i,j)\;=\;\int\!\!\int
  \Cov_X\bigl(M_{w,L,\cdot}(\rho),\,M_{w,L,\cdot}(\rho')\bigr)
  \,d\mu_{ij}(\rho,\rho') ,
\]
the covariance being taken in the window position $x$ under $\mathbb
P_X$, so that $K_{\mathrm{spec}}$ is the window pullback of the joint
zero-correlation measure and of nothing else.
\emph{(ii)} [strong converse layer] no fourth covariance source survives
at leading order, that is, the residual after subtracting the three
displayed terms is $o(V)$ for every admissible pair.  This clause is
separately falsifiable from (i), and a mixed term representable by
neither connected local densities nor zero correlations refutes it alone.
\emph{(iii)} [crossover scale, derived] the crossover length $L_*(X)$ is
defined by equality of the Hilbert--Schmidt norms of the two $m\times m$
matrix kernels,
\[
  \bigl\|K_{\mathrm{HL}}(L)\bigr\|_{\mathrm{HS}}
  \;=\;\bigl\|K_{\mathrm{spec}}(L)\bigr\|_{\mathrm{HS}},
  \qquad\|A\|_{\mathrm{HS}}^{2}=\sum_{i,j}|A_{ij}|^{2},
\]
this being the sense of ``equal mass'' used throughout, and $L_*(X)$ is
therefore determined by (i) and not by a fitted exponent.
\end{conjecture}

Local congruence correlations and global zero oscillations are not
competing metaphors but two expansions of the same von Mangoldt field,
and the coincident shifts must be separated first because they belong to
neither description.  A proof would join short-interval prime statistics,
tuple singular series and zero correlations into one calculus, and would
say when a random-sieve model suffices, when spectral information is
unavoidable, and which residual statistic would signal genuinely new
arithmetic structure.  Explicit-formula methods and Hardy--Littlewood
tuple expansions are the two established halves, developed for prime
counts by Montgomery and Soundararajan and for the singular series by
Pintz, and the claim distinctive here is that they are exactly the two
pieces of one complete decomposition for arbitrary motifs.

\emph{Evidence.}  The twin-pair Fano factor fell from $0.83$ to $0.71$ as
the window mean rose from $2.49$ to $93.31$ while the skewness moved
toward zero, which is consistent with persistent sub-Poisson rigidity
across scale but cannot separate the connected local kernel from a zero
contribution at this range, and the windows again overlap at step $L/4$
so that roughly a quarter of the nominal samples are effective and no
error bars exist.

\begin{conjecture}[Spectrally conditioned extremal index for prime
motifs]\label{c4}
Fix an admissible motif $H$ and the window-position law (W) on
$\Omega_X=[X,2X]$ with window length $L$.  Let $N_H(x)$ be the point
process of occurrences of $H$ in $[x,x+L]$, transformed by the exact
cumulative Hardy--Littlewood hazard $t\mapsto\SG(H)\int_2^{t}(\log
s)^{-|H|}ds$ so that it has unit intensity.  Fix a cutoff wavelength
$L_{\mathrm{spec}}$ with $L_{\mathrm{spec}}/L\to\infty$, the symbol
$\Lambda$ being reserved for the von Mangoldt function, let
$\Psi_{L_{\mathrm{spec}}}(x)$ be the restriction to $[x,x+L]$ of the
truncated explicit-formula sum over the zeros of the relevant
$L$-functions with $|\Im\rho|\le X/L_{\mathrm{spec}}$, and let $\mathcal
F_{L_{\mathrm{spec}}}=\sigma\bigl(\Psi_{L_{\mathrm{spec}}}\bigr)$ be the
sub-sigma-field of the Borel field of $\Omega_X$ that it generates.

The \emph{admissible self-overlap graph} $O(H)$ has for vertex set
\[
  V(H)\;=\;\bigl\{t\in\mathbb Z\setminus\{0\}:\ (H+t)\cap
  H\neq\emptyset\ \text{and}\ H\cup(H+t)\ \text{admissible}\bigr\},
\]
a finite set contained in $(H-H)\setminus\{0\}$, and two vertices $t\neq
t'$ are joined by an edge exactly when the three translates share a
further coincidence and the triple union is admissible, that is, when
\[
  (H+t)\cap(H+t')\neq\emptyset,\quad\text{equivalently}\quad
  t-t'\in(H-H)\setminus\{0\},
\]
and the union $H\cup(H+t)\cup(H+t')$ is admissible.
For a set $S\subseteq V(H)$ write $\SG(S)=\SG\bigl(H\cup\bigcup_{t\in
S}(H+t)\bigr)$ for the Bateman--Horn density of the corresponding union
pattern, so that $\SG(\emptyset)=\SG(H)$.  Let $\mathcal K(H)$ be the set
of connected components of $O(H)$ together with the empty component
$\emptyset$, which records an isolated occurrence.  The
\emph{visited-component law} is the probability measure on $\mathcal
K(H)$ given by
\[
  \mathbb P\bigl[\mathcal C\bigr]\;=\;
  \frac{\SG\bigl(V(\mathcal C)\bigr)}
  {\sum_{\mathcal C'\in\mathcal K(H)}\SG\bigl(V(\mathcal C')\bigr)} ,
\]
the weight of a component being proportional to the Bateman--Horn density
of the constellation it realizes, and the size of the visited component
is $|V(\mathcal C)|+1$, the base occurrence included, so that the empty
component has size one.  Then:
\emph{(i)} the conditional law of the cluster starts of $N_H$ given
$\mathcal F_{L_{\mathrm{spec}}}$ converges, as $X\to\infty$ and in
probability under $\mathbb P_X$, to that of a homogeneous Poisson
process.
\emph{(ii)} the within-cluster law is the visited-component law just
defined, so the conditional extremal index is the reciprocal mean
component size
\[
  \theta(H)\;=\;\Bigl(\sum_{\mathcal C\in\mathcal K(H)}\mathbb P[\mathcal
  C]\bigl(|V(\mathcal C)|+1\bigr)\Bigr)^{-1}
\]
and involves no fitted parameter.
\emph{(iii)} after the conditioning is removed, that is, after averaging
over $\mathcal F_{L_{\mathrm{spec}}}$ under $\mathbb P_X$, the limit is
the Cox cluster process directed by the exponential of the limiting
low-mode field.
\emph{(iv)} [functoriality] if two motifs $H$ and $H'$ have isomorphic
admissible self-overlap graphs, then $\theta(H)=\theta(H')$ once each
process has been put on the unit-intensity clock by its own cumulative
Hardy--Littlewood hazard, which is the local-density normalization meant
here: the transform divides by $\SG(H)(\log t)^{-|H|}$, so the extremal
index compares cluster structure and not the differing occurrence rates
of the two motifs.  Motifs with empty $O(H)$, that is with
$V(H)=\emptyset$, admit only the empty component and have
$\theta(H)=1$.
All four clauses are assertions about the probability space
$(\Omega_X,\mathbb P_X)$.  The spectral field is deterministic, what is
random is the window position, and $\mathcal F_{L_{\mathrm{spec}}}$ is
the sigma-field the window position generates through the
long-wavelength component.
\end{conjecture}

Short overlaps create genuine clusters while zeros and other long modes
modulate the local intensity, and conditioning separates the two
mechanisms instead of absorbing both into a fitted location shift of a
Gumbel law.  The statement makes extreme prime-pattern gaps a probe of
local arithmetic geometry and of the low-frequency prime field at once,
predicting null cases with $\theta(H)=1$, exact changes on adjoining a
shift, and a principled source for apparent non-Gumbel mixtures.  Generic
extreme-value limits in the Gumbel class are classical and do not
identify the arithmetic source of clustering, so the claimed content is
the conditional cluster process together with the identification of its
within-cluster law with the admissible self-overlap graph.

\emph{Evidence.}  For $17{,}937$ twin-pair gaps in the upper part of the
tested range the mean-scaled hazard had coefficient of variation $0.904$
and Kolmogorov--Smirnov distance $0.060$ from $\Exp(1)$, and since the one
per cent critical value at that sample size is about $0.012$, the
unconditioned hazard law is decisively non-exponential.  The direction
matters.  A coefficient of variation of $0.904$ is below one and is
therefore underdispersion, whereas every mixture of exponentials, and in
particular every Cox law of the form asserted in clause (iii), has
coefficient of variation at least one.  The measurement therefore argues
against the bare exponential and against an unconditioned Cox reading
alike, and is informative about rigidity at this range rather than about
the mixture mechanism.  No spectral conditioning is available here, so
clauses (i) and (ii) are untested.

\begin{conjecture}[Completeness of arithmetic wavelets]\label{c5}
Fix a finite family of admissible motifs $H_1,\dots,H_m$, the law (W),
and the incidence tensors $A^{(j)}_r$ of Section~\ref{sec:notation},
together with the renormalized disjoint kernel $K_{\mathrm{HL}}$, which
is examined last.  Write $\mathcal A=\bigcap_{e\ge1}\ker A^{(1,e)}$ for
the joint kernel of the second-order coincidence matrices.  Then:
\emph{(i)} [constructive layer] if $c\in\mathbb R^{m}$ satisfies
$A^{(1,e)}c=0$ for every $e\ge1$ and $c^{\mathsf T}K_{\mathrm{HL}}(L)c
\neq0$, then $\Var_XY_c$ begins exactly at the renormalized disjoint
kernel, with the scale and constant predicted by Conjectures~\ref{c2}
and~\ref{c3}, and if $c\notin\mathcal A$ then $\Var_XY_c$ begins at
$L(\log X)^{e}c^{\mathsf T}A^{(1,e)}c$ for the largest $e$ with
$A^{(1,e)}c\neq0$.
\emph{(ii)} [strong converse layer] conversely, every sequence of
coefficient vectors whose variance is smaller than the generic
coincidence scale $L(\log X)^{e_{\max}}$ by a fixed power of $L$ or by a
fixed power of $\log X$ converges in direction to $\mathcal A$.
\emph{(iii)} the space $\mathcal A$ is therefore the first canonical
arithmetic wavelet space attached to the family, and the second-order
filtration $\mathbb R^{m}\supseteq\mathcal A$ has depth one, in
accordance with the two-level overlap poset of a pair.  Its deeper terms
are not visible in the variance and are cut out by the higher-order
tensors, the space of level $(r,s)$ being $\bigcap_{j\le s}\ker
A^{(j)}_r$ for $r\ge3$, so that the arithmetic wavelet spaces of the
family form a decreasing filtration indexed by the cumulant order as well
as by the codimension.
Clause (ii) is separately falsifiable from (i): a rigid coefficient
sequence generated by a global spectral cancellation not represented in
the finite overlap tensors refutes (ii) while leaving (i) and (iii)
untouched.
\end{conjecture}

Cancellation of overlap diagrams is the arithmetic analogue of vanishing
moments in wavelet theory, and the converse layer asserts that there is
no other route to rigidity, so the finite overlap and spectral tensors
exhaust the ways a balanced prime statistic can be made quiet.  The
practical consequence is that low-variance combinations can be
constructed on purpose to expose deeper correlation strata, isolating
zero terms, distinguishing competing tuple conjectures and giving
high-sensitivity numerical tests.  Balanced combinations of
prime-counting functions are common in the literature, in the prime-race
tradition surveyed by Martin and coauthors among others, but a converse
classification of all unusually rigid finite motif statistics by overlap
annihilation is the claimed content.

\emph{Evidence.}  The empirically smallest covariance mode of the four
pair fields carried $0.58$ to $0.59$ of the mean coordinate variance at
every window and its coefficient vector reached absolute alignment
$0.987$ between the two largest windows while remaining unstable at small
windows, but that mode is selected by the same data that measure its
variance, the windows overlap at step $L/4$ so that about a quarter of
the nominal samples are effective, and no error bar accompanies the
alignment.

\begin{question}\label{q:kernel}
For four pair motifs, is the joint kernel $\mathcal
A=\bigcap_{e\ge1}\ker A^{(1,e)}$ of Conjecture~\ref{c5}
nonzero over $\mathbb Q$, and does the rational vector spanning it agree
in direction with the empirically smallest covariance mode?  This is an
open question and is not counted among the twenty-five conjectures.
\end{question}

\section{Programme II. Arithmetic occupancy and local-global
decoupling}\label{sec:II}

Least primes in residue classes form a field carrying both an additive
exclusion geometry and a multiplicative spectrum, and odd class groups
are tested against a growing set of local conditions.  Conjectures
\ref{c6} to \ref{c9} live on the modulus-class law (M), and
Conjecture~\ref{c10} on the uniform discriminant ensemble declared in its
statement.

\begin{conjecture}[First-order arithmetic occupancy kernel]\label{c6}
Work on the modulus-class law (M) and put
$T_q(a)=\Li(p(q,a))/(q-1)$.  Let $S_Q(t)=\mathbb P_Q\bigl[T_q(a)>t\bigr]$
be the survival function of the ensemble.  Then, uniformly for $t$ in
compact subsets of $[0,\infty)$,
\[
  S_Q(t)\;=\;\e^{-t}\Bigl[1-\frac{A(t)}{\log Q}+o\Bigl(\frac1{\log
  Q}\Bigr)\Bigr]\qquad(Q\to\infty),
\]
where $A$ is the canonical function defined as follows.  Expanding the
empty-class probability $\mathbb P_Q[N_a(y)=0]$, with
$N_a(y)=\#\{p\le y:p\equiv a\ (q)\}$, by inclusion and exclusion in
connected prime correlations produces at second order the connected
two-prime exclusion kernel $\SG(\{0,kq\})-1$, $k\ge1$, of
Section~\ref{sec:notation}, and $A(t)$ is the coefficient of $1/\log Q$
obtained by applying the first-arrival transform $\mathcal F$ to that
kernel and integrating to $\Li$-time $t$.  In particular $A$ carries no
fitted parameter and $A(0)=0$.
[stated programme] The closed evaluation of $A$ from the numbers
$\SG(\{0,kq\})$, the passage from the discrete kernel indexed by $k$ to
the bounded kernel on $\mathbb R_{\ge0}$ on which the first-arrival
transform $\mathcal F$ of Section~\ref{sec:notation} acts, which is the
change of variable to $\Li$-time, and the proof that all higher connected
corrections are $O(1/\log^2Q)$ after averaging $q$ over the dyadic block,
are the derivation route registered with this conjecture.
\end{conjecture}

The independent collector gives the bare exponential, but nearby primes
cannot occupy residue classes independently, since divisibility and
prime-pair constraints impose a short-range exclusion process, and
summing that exclusion produces the first arithmetic deformation of the
survival law.  A proof would supply the first quantitative stochastic
theory of the whole least-prime field rather than an order of magnitude
for the extreme class, and it feeds Conjectures~\ref{c7} to~\ref{c9}
directly and generalizes to Chebotarev occupancy.  Li, Pratt and Shakan
formulate least primes as an arithmetic coupon collector and obtain the
extremal lower bound in that language, and the content claimed here is
the complete first-order survival deformation with its kernel derived
from connected prime correlations.

\emph{Evidence.}  Across prime moduli from $101$ to $999$ the scaled mean
speedup $(1-\mathbb E_a[T_q])\log q$ rose from $1.27$ to $1.59$ and the
cover time sat $2.12$ to $2.37$ below the independent harmonic-mean
value, so the sign and the persistence are inconsistent with an
exchangeable collector at first order, but $\log q$ varies by less than a
factor of $1.5$ across the whole range, which cannot separate $1/\log Q$
from a $\log\log Q/\log Q$ correction.

\begin{conjecture}[Dual additive-multiplicative spectral geometry of
least primes]\label{c7}
Work on the modulus-class law (M), with $\widetilde T_q$,
$S^{\mathrm{add}}_q$ and $S^{\mathrm{mult}}_q$ as defined in
Section~\ref{sec:notation}.  Then:
\emph{(i)} [additive half] for every fixed $\alpha\in(0,1)$,
\[
  \mathbb E_Q\bigl[S^{\mathrm{add}}_q(\lfloor\alpha q\rfloor)\bigr]
  \;=\;\widehat K_{\mathrm{add}}(\alpha)
  \;+\;o\Bigl(\frac1{\log Q}\Bigr),
\]
where $\widehat K_{\mathrm{add}}$ is the Fourier transform of $\mathcal
F$ applied to the connected two-prime exclusion kernel, an explicit
function carrying no fitted parameter.  [stated programme] The closed
evaluation of $\widehat K_{\mathrm{add}}$ from the numbers
$\SG(\{0,kq\})$ through the transform $\mathcal F$ of
Section~\ref{sec:notation} is the derivation route registered here.
\emph{(ii)} [multiplicative half] characters of exact order $k$ modulo
$q$ exist only when $k$ divides $q-1$, so the average is taken over the
subensemble of (M) carried by the moduli $q\equiv1\pmod k$, renormalized
to a probability measure on that subensemble, which is nonempty and of
positive mass for every fixed $k$ by Dirichlet's theorem.  For every
fixed order $k\ge2$ the average of $S^{\mathrm{mult}}_q(\chi)$ over the
Dirichlet characters of exact order $k$ modulo $q$ equals
$\lambda_k+o(1/\log Q)$ under that conditioned expectation, where
$\lambda_k$ is the explicit-formula contribution of the zeros of the
corresponding $L$-functions to the first-arrival functional.  [stated
programme] The closed evaluation of $\lambda_k$ from those zeros through
$\mathcal F$ is the derivation route registered here.
\emph{(iii)} both halves carry mass of the same order, so the field is
diagonal in neither the additive nor the multiplicative basis, its
covariance being the sum of two canonical operators of low complexity,
one a convolution on the additive group of $\mathbb Z/q\mathbb Z$ and one
diagonal in the Dirichlet characters.
\emph{(iv)} [strong converse layer] the two halves exhaust the field, in
that $\mathbb E_Q\bigl[\|\widetilde T_q\|_2^2\bigr]$ equals the total
mass supplied by (i) and (ii) up to $o(1/\log Q)$, so no third invariant
is required and the mixed additive-multiplicative correlations vanish
after averaging $q$ over the dyadic block.  Clause (iv) is separately
falsifiable from (i) to (iii): a persistent mixed component refutes (iv)
alone and forces a third operator rather than a correction to the first
two.
\end{conjecture}

Prime differences see the additive group of $\mathbb Z/q\mathbb Z$
whereas the distribution among reduced classes sees its character group,
and the first-arrival functional is nonlinear enough to retain both
sources at first order without mixing them.  A proof would give a genuine
covariance theory for the entire least-prime landscape and would explain
why purely independent, purely additive and purely prime-race models are
each incomplete, and it supplies the environment conditioned on in
Conjecture~\ref{c8} and the background against which the spike of
Conjecture~\ref{c9} is defined.  Prime races are expanded in
multiplicative characters in the tradition surveyed by Martin and
coauthors, and effective Chebotarev inputs of the kind refined by Thorner
and Zaman govern the character side, while the exclusion between nearby
primes is an additive phenomenon in the sense of Li, Pratt and Shakan, so
the claim is that the least-prime field contains both geometries and that
they exhaust its first-order covariance.

\emph{Evidence.}  The $q$-point transform on the full additive group,
including the zero class, at which the recorded value is
$p(q,0)=q$ and therefore $T_q(0)=\Li(q)/(q-1)$, gave low-frequency to
high-frequency power ratios of $1.13$ to $1.42$ across modulus bands,
which rejects a flat
spectrum and supports the additive half, while the multiplicative
covariance was not estimated at all, so clause (i) is tested on one of
its two operators and clause (iii) is untested.

\begin{conjecture}[Conditioned Cox--Poisson law for the last residue
classes]\label{c8}
Work on the modulus-class law (M), so that the modulus $q$ is uniform on
the primes of $[Q,2Q]$ and the centred arrival field
$a\mapsto\widetilde T_q(a)$ is a random element of $\mathbb R^{q-1}$.
Say that a mode of $\widetilde T_q$ is of \emph{terminal scale} if its
contribution to the variance of the field on the last-arrival window,
that is to $\Var$ of $a\mapsto T_q(a)-\log(q-1)$, stays bounded away from
$0$ and from $\infty$ as $Q\to\infty$.  Conjecture~\ref{c7} lists them
explicitly: they are the additive frequencies $\xi$ with $|\xi|\le\xi_0$,
where $\xi_0$ is the fixed cutoff beyond which $\widehat
K_{\mathrm{add}}(\xi/q)$ falls below a fixed positive threshold, together
with the Dirichlet characters of exact order $k\le k_0$, where $k_0$ is
the fixed cutoff beyond which $\lambda_k$ falls below the same threshold.
The list is finite and does not depend on $q$.  Let $Z_q$ be the
projection of $\widetilde T_q$ onto those modes, a random vector on
$\Omega_Q$, and let $\mathcal G_Q=\sigma(Z_q)$ be the sub-sigma-field it
generates.  Let $c_6(q)$ be the deterministic centring supplied by
Conjecture~\ref{c6}, namely the value that turns the survival law of that
conjecture into a unit-rate exponential tail at the last-arrival level,
\[
  c_6(q)\;=\;\log\Bigl(1-\frac{A\bigl(\log(q-1)\bigr)}{\log q}\Bigr)
  \;=\;-\frac{A\bigl(\log(q-1)\bigr)}{\log q}
  +O\Bigl(\frac1{\log^{2}q}\Bigr),
\]
the extension of $A$ from compact arguments to arguments of size $\log q$
being part of the stated programme registered with
Conjecture~\ref{c6}.  Let
\[
  \Xi_q\;=\;\sum_{a\in(\mathbb Z/q\mathbb Z)^{\times}}
  \delta_{\,T_q(a)-\log(q-1)-c_6(q)}
\]
be the point process of last-arrival excesses.  Then:
\emph{(i)} conditionally on $\mathcal G_Q$, $\Xi_q$ converges in law
under $\mathbb P_Q$ to a homogeneous Poisson process of intensity
$\theta_{\mathrm{occ}}\e^{-x}\,dx$.
\emph{(ii)} unconditionally, $\Xi_q$ converges to the Cox process
directed by $\theta_{\mathrm{occ}}\exp(-x+W)\,dx$, where $W$ is the
limiting low-mode field of Conjecture~\ref{c7}.
\emph{(iii)} the constant $\theta_{\mathrm{occ}}$ is the cluster
functional of the short-range occupancy kernel of Conjecture~\ref{c6} and
is therefore derived rather than fitted.
Every clause is a statement about the probability space
$(\Omega_Q,\mathbb P_Q)$.  The arrival field of a single modulus is
deterministic, what is random is the modulus, and $\mathcal G_Q$ is the
sigma-field generated by the finitely many mode coordinates of the
sampled field.
\end{conjecture}

Rare unfilled classes are approximately independent only after the slowly
varying arithmetic environment is held fixed, since the low modes
randomize the local intensity while the short exclusion kernel changes
the extremal index, and the two effects are separated by exactly the
conditioning declared above.  The statement predicts the joint law of the
last several classes, their spacings, multiplicities and character
profiles, and it is a sharp diagnostic, because pure Poisson, Cox--Poisson
and clustered Poisson models make different predictions for the
probability of two exceedances.  The classical coupon collector has a
Poisson terminal process, and the point of departure is that the
arithmetic field of Conjecture~\ref{c7} has persistent low modes, so an
unconditioned Poisson law is too rigid.

\emph{Evidence.}  The tested moduli are too small for the process-level
claim, since the centred exceedance counts above $0$, $1$ and $2$ were
usually sparse or zero across the prime moduli below $1000$, so clauses
(i) and (ii) are untested and only the stable negative cover-time shift,
which supports a nontrivial centring in clause (iii), was observed.

\begin{conjecture}[Rank-one exceptional-zero deformation of
occupancy]\label{c9}
Fix $\varepsilon_0=10^{-2}$ and let $\varepsilon\le\varepsilon_0$ be
fixed.  Work on the modulus-class law (M) restricted to those $q$
admitting a primitive real character $\chi_q$ whose $L$-function has a
real zero $\beta$ with $(1-\beta)\log q\le\varepsilon$, the limits being
taken along any sequence of $Q$ for which this subensemble is nonempty
and carries positive mass.  The statement is made for every fixed
$\varepsilon\le\varepsilon_0$, the value of $\varepsilon_0$ serving only
to fix the meaning of ``close to $1$'' and no clause depending on it
otherwise.  Under the generalized Riemann hypothesis the subensemble is
empty for every $Q$, and all three clauses are then instanceless rather
than false, which is why the conjecture is stated conditionally on the
subensemble carrying positive mass.  Then, after the ordinary occupancy
kernel of Conjecture~\ref{c6} is removed:
\emph{(i)} the centred first-arrival field has a deterministic component
along $\chi_q$ whose amplitude is the first-arrival transform $\mathcal
F$ applied to the exceptional explicit-formula term $x^{\beta}/\beta$,
an explicit function of $\beta$ and $q$ with no fitted parameter.
\emph{(ii)} all modes orthogonal to $\chi_q$ obey Conjecture~\ref{c8} to
first order.
\emph{(iii)} [strong converse layer] no other first-order deformation
survives, so the effect of an exceptional zero on the nonlinear
first-arrival field is asymptotically of rank one.  Clause (iii) is
separately falsifiable from (i) and (ii), and higher harmonics of
$\chi_q$ generated by the nonlinearity would refute it alone.
\end{conjecture}

The explicit formula inserts a single exceptional term proportional to
$\chi_q(a)$, and linearizing the first-hit functional transfers that term
to a rank-one perturbation of the field, in the manner of a spiked random
matrix or a spiked random field.  If true the statement makes least-prime
occupancy a sensitive diagnostic for exceptional zeros and separates
ordinary arithmetic dependence from zero-driven bias precisely, and if
false the failure locates the nonlinearity that generates the missing
harmonics.  Effective Chebotarev and comparative prime number theory, in
the forms developed by Thorner and Zaman and surveyed by Martin and
coauthors, establish that real zeros create character biases in linear
counts, and the new claim is that their effect on the nonlinear
first-arrival field is of rank one.

\emph{Evidence.}  No support is available at the accessible scale, the
Spearman correlation between the absolute quadratic-character projection
and an inverse $L(1,\chi)$ proxy being $0.066$ over the prime moduli
below $1000$, a null measurement rather than evidence against the
conjecture, since no exceptional zero occurs in that range.

\begin{conjecture}[Growing local-condition independence for odd class
groups]\label{c10}
Fix an odd prime $\ell$.  Sample $D$ uniformly from the negative
fundamental discriminants with $|D|\le Y$, and condition on the complete
two-primary genus data of $D$ in the sense of
Section~\ref{sec:notation}.  Let $B=B(Y)$ and let
$\Sigma_B(D)=\bigl((D/p)\bigr)_{p\le B}$ be the vector of splitting
signs.  The single information measure used throughout is the number of
imposed local bits $\pi(B)$, compared against the logarithmic information
content $\log Y$ of the discriminant.  Then:
\emph{(i)} whenever $\pi\bigl(B(Y)\bigr)=o(\log Y)$, the conditional law
of $\mathrm{Cl}(D)[\ell^{\infty}]$ given the genus data and given
$\Sigma_B(D)$ converges in total variation, as $Y\to\infty$ and uniformly
over the attainable sign vectors, to the Cohen--Lenstra law for the odd
part conditioned on the genus data alone.
\emph{(ii)} the total-variation error is governed explicitly by
$\pi(B)/\log Y$ and tends to zero whenever $\pi\bigl(B(Y)\bigr)=o(\log
Y)$, which is the same range as in (i).  No sharpness is claimed at the
boundary $\pi(B)\asymp\log Y$, where the conjecture makes no assertion.
\emph{(iii)} the assertion is made for the full isomorphism type of
$\mathrm{Cl}(D)[\ell^{\infty}]$ and not only for the divisibility of
$h(D)$ by $\ell$, the latter being a coarse projection of the former.
\end{conjecture}

Genus theory is the known global coupling between local data and the
class group, and once it is projected out, Frobenius data at bounded
primes should consume only local entropy and should not bias the odd
unramified extension group until the number of imposed conditions
approaches the logarithmic information content of the discriminant.  A
proof would be a quantitative local-global decoupling principle joining
Chebotarev data to Cohen--Lenstra statistics, would say how much local
information a class group can absorb, and has direct analogues for Selmer
groups and ray class groups.  Wood establishes Cohen--Lenstra behaviour
under fixed local conditions, Klagsbrun proves Davenport--Heilbronn
theorems for quotients of class groups, Wang and Wood compute the moments
that interpret the heuristics and Ellenberg surveys the field, and the
content claimed here is uniformity over a growing set of Frobenius
conditions after the genus obstruction is removed.

\emph{Evidence.}  For the $4560$ negative fundamental discriminants with
$|D|\le15{,}000$, conditioning on a single splitting sign at
$p=3,5,7,11,13$ moved the probability that $3$ divides $h(D)$ by at most
$0.012$ and that $5$ divides $h(D)$ by at most $0.026$, but the test
records divisibility of the class number and not the isomorphism type
required by clause (iii), and the multi-cell variation is dominated by
small samples.

\section{Programme III. Finite logarithms, Kummer rank and algebraic
tori}\label{sec:III}

Finite logarithms live on relation tori, saturated rank and $p$-adic
depth control the exceptional strata, and the same principle extends
functorially from the multiplicative group to arbitrary algebraic tori.
Every statement in this programme is asserted only for saturated groups
and only after the Kummer degrees of Section~\ref{sec:notation} have been
computed.

\begin{conjecture}[Kummer--Haar law for Fermat-quotient
homomorphisms]\label{c11}
Let $\Gamma\le\mathbb Q^{\times}$ be finitely generated and saturated of
rank $r$, fix an ordered generating tuple $g=(g_1,\dots,g_s)$ with
relation lattice $R(g)$ and ambient torus $(\mathbb R/\mathbb Z)^{s}$ as
in Section~\ref{sec:notation}, and let $p$ range over the primes dividing
no numerator or denominator of a generator.  Then:
\emph{(i)} [elementary] $q_p(ab)\equiv q_p(a)+q_p(b)\pmod p$, so
$q_p\colon\Gamma\to\mathbb F_p$ is a homomorphism and the ambient vector
$\mathbf q_p(g)/p\in(\mathbb R/\mathbb Z)^{s}$ lies in the subtorus
$\mathbb T_\Gamma$ cut out by the relation lattice $R(g)$, a subtorus of
dimension exactly $r$.
\emph{(ii)} as $p$ varies through any fixed compatible Chebotarev class,
defined by prescribed splitting behaviour in a fixed cyclotomic-Kummer
field, $\mathbf q_p(g)/p$ equidistributes with respect to Haar measure on
$\mathbb T_\Gamma$.
\emph{(iii)} [strong converse layer] the support is exactly
$\mathbb T_\Gamma$ and no further constraint holds: for every nonrelation
character $m\in\mathbb Z^{s}$ of the tuple, that is every $m$ outside the
saturation $R(g)^{\mathrm{sat}}$, the mean of $\exp\bigl(2\pi i\langle
m,\mathbf q_p(g)\rangle/p\bigr)$ over the class tends to zero.  The
exclusion of $R(g)^{\mathrm{sat}}$ is necessary and sufficient, since on
that sublattice the exponential is identically $1$ at all but finitely
many $p$, whereas every $m$ outside it restricts nontrivially to $\mathbb
T_\Gamma$.  Clause (iii) is separately falsifiable from
(ii), a persistent Fourier coefficient at a nonrelation frequency
refuting it alone.
\emph{(iv)} [invariance, with the saturation caveat in force] the law
depends on $\Gamma$ only through $\Gamma^{\mathrm{sat}}$, so two
finitely generated groups with the same saturation give the same limit
law and the controlling invariant is the saturated rank $r$ and not the
number of displayed generators.  For a group that is not saturated,
clause (iv) is false as stated with $r$ read as the rank of the displayed
generators, which is why saturation is part of the hypothesis rather than
a proviso.
\end{conjecture}

The Eisenstein--Lerch identity makes $q_p$ a finite logarithm, Kummer
theory identifies the exact relation lattice that its values must
respect, and the conjecture asserts that no constraint beyond that
lattice survives, so the limit is Haar on the quotient torus.  A proof
would supply a single structural law covering ordinary Fermat quotients,
simultaneous near-Wieferich events, restrictions to arithmetic
progressions and mod-$p$ regulators, with rank replacing the number of
bases as the controlling invariant.  Shparlinski gives exponential-sum
and value-set bounds for Fermat quotients and the horizontal statistics
of Ostafe and Shparlinski describe the fixed-$p$ picture, while the
mod-$p$ Leopoldt viewpoint of Boeckle, Guiraud, Kalyanswamy and Khare
connects exceptional vanishing to regulators, and the full rank-torus law
above is the claimed content.

\emph{Evidence.}  For the $22{,}041$ primes $p\le250{,}000$ the exact
relations $q_p(6)=q_p(2)+q_p(3)$ and $q_p(12)=2q_p(2)+q_p(3)$ held to
within $3\times10^{-16}$, which checks clause (i) and the implementation
rather than the conjectural clauses, and the nontrivial one- and
two-dimensional Fourier coefficients for bases $2,3,5$ stayed below
$0.011$, of the square-root size $1/\sqrt{22041}\approx0.0067$ expected
for that sample, so clause (iii) is tested only at fixed frequencies.

\begin{conjecture}[Square-root large sieve for finite
logarithms]\label{c12}
For every saturated $\Gamma\le\mathbb Q^{\times}$ of rank $r$ and every
ordered generating tuple $g=(g_1,\dots,g_s)$ of $\Gamma$ there exists
$\delta(\Gamma)>0$ such that, for every nonrelation character
$m\in\mathbb Z^{s}$ of the tuple in the sense of
Section~\ref{sec:notation}, that is every $m$ outside the saturated
relation lattice $R(g)^{\mathrm{sat}}$, of height $\max_i|m_i|$ at most
$X^{\delta(\Gamma)}$,
\[
  \Bigl|\sum_{X<p\le2X}\e\Bigl(\frac{\langle m,\mathbf
  q_p(g)\rangle}{p}\Bigr)\Bigr|
  \;\ll_{\Gamma,\varepsilon}\;\pi(X)^{1/2}X^{\varepsilon},
\]
uniformly after conditioning on fixed compatible cyclotomic and Kummer
data, and the same bound holds for any fixed finite set of joint
frequencies, the square-root scale being set by the number of prime
samples and not by the ambient coordinate count.  The growing height
range is the essential content: the corresponding statement for a single
fixed frequency is weak equidistribution, whereas the range
$X^{\delta(\Gamma)}$ is what controls targets of shrinking measure.
\end{conjecture}

Once the exact algebraic relations are quotiented out, finite logarithms
should oscillate as genuinely independent phases, and the conjecture
demands exactly the square-root cancellation that independence predicts,
neither more, which would be stronger than the model supports, nor less.
A proof would be a large sieve theorem for $p$-adic logarithms, would
make Conjecture~\ref{c11} quantitative, and would yield the
shrinking-target estimates that Conjectures~\ref{c13} and~\ref{c14}
require and the toric analogues of Conjecture~\ref{c15}.  The existing
Fermat-quotient bounds of Shparlinski concern other averaging regimes and
shorter frequency ranges, so the prime-varying, rank-sensitive and
height-growing estimate above is the claimed strengthening.

\emph{Evidence.}  At frequencies $1,2,3,5,8,13$ the normalized Fourier
magnitudes for bases $2,3,5$ lay between $0.001$ and $0.010$ on about
$22{,}000$ primes, which is of the square-root size for that sample, but
every frequency tested is fixed, so the computation bears on none of the
growing-height content that distinguishes this statement from weak
equidistribution.

\begin{conjecture}[Saturated-rank threshold for simultaneous Wieferich
primes]\label{c13}
Let $\Gamma\le\mathbb Q^{\times}$ be finitely generated, let
$\Gamma^{\mathrm{sat}}$ be its saturation and $r=\rank\Gamma$, and let
$\kappa(\Gamma;p)=p^{r}/\Delta_\Gamma(p)\ge1$ be the local Kummer factor
of Section~\ref{sec:notation}.  All clauses are asserted for
$\Gamma=\Gamma^{\mathrm{sat}}$ and only after the Kummer degrees
$\Delta_\Gamma(p^{k})$ have been computed.  Put
$W(\Gamma;X)=\#\{p\le X:\ q_p\ \text{vanishes on}\ \Gamma\}$.  Then:
\emph{(i)} if $r=1$ then $W(\Gamma;X)\sim\kappa(\Gamma)\log\log X$, with
\[
  \kappa(\Gamma)\;=\;\lim_{X\to\infty}\frac1{\log\log X}
  \sum_{p\le X}\frac{\kappa(\Gamma;p)}{p},
\]
a compatible Kummer density determined by the degrees
$\Delta_\Gamma(p)$ and equal to $1$ when $\Gamma$ is Kummer-generic at
almost all $p$.
\emph{(ii)} if $r\ge2$ then $\sum_p\kappa(\Gamma;p)p^{-r}$ converges and
only finitely many such primes exist, the expected tail beyond $Y$ being
$\sum_{p>Y}\kappa(\Gamma;p)p^{-r}$.
\emph{(iii)} [invariance, with the saturation and compatibility caveats
in force] adjoining to $\Gamma$ a generator lying in
$\Gamma\otimes\mathbb Q$ changes no term of (i) or (ii), because the
saturation and hence $r$ and every $\kappa(\Gamma;p)$ are unchanged,
whereas adjoining a generator independent of $\Gamma$ multiplies the
local term at $p$ by $1/p$ exactly at those $p$ at which the enlarged
group remains Kummer-generic, the factor being
$\kappa(\Gamma';p)/\bigl(p\,\kappa(\Gamma;p)\bigr)$ at the finitely many
entangled primes.  Both halves of (iii) fail for a group that is not
saturated, since the displayed rank is then not the exponent, and the
second half fails for a Kummer-entangled extension, since the degree
$\Delta$ is then not multiplied by $p$, so saturation and Kummer
compatibility are hypotheses of the statement and not annotations to it.
\end{conjecture}

Haar measure on an $r$-dimensional finite-logarithm torus assigns mass
$p^{-r}$ to the origin, so the divergence or convergence of the
Borel--Cantelli sum is decided by the codimension alone, and Kummer
splitting changes the local constant without touching that codimension.
The statement unifies ordinary, simultaneous and relation-dependent
Wieferich questions and gives an algebraic explanation of the threshold,
and it is the direct arithmetic test of Conjecture~\ref{c11}.  The
one-base heuristic with probability $1/p$ is classical folklore, the
mod-$p$ Leopoldt framework of Boeckle, Guiraud, Kalyanswamy and Khare
supplies the broader regulator setting, and Silverman derives
non-Wieferich statements from the $abc$ conjecture, with progressions by
Graves and Murty, so the sharpening
claimed here is that the saturated Kummer rank, and not the displayed
number of bases, is the exact exponent.

\emph{Evidence.}  For symmetric near-zero boxes of half-width $0.02$ the rank
one, two and three counts were $922$, $33$ and $1$ against Haar-volume
predictions of about $882$, $35$ and $1.4$, and exact vanishing occurred
for individual bases but for no pair of bases below $250{,}000$, which is
compatible with clause (ii) at a range far too small to test a
convergent tail.

\begin{conjecture}[Rank-depth Poisson law for higher Wieferich
strata]\label{c14}
Fix $r\ge1$ and $\theta<1$, and sample an ordered $r$-tuple of
multiplicatively independent integers of height at most $H$ uniformly
from the admissible tuples, admissibility excluding the
Kummer-incompatible tuples of Section~\ref{sec:notation}.  Let
$\Gamma(H)$ be the saturation of the sampled tuple.  Mark the pair
$(p,k)$, for $p\le H^{\theta}$ and $k\ge1$, when every generator $a$
satisfies $a^{p-1}\equiv1\pmod{p^{k+1}}$, and put
\[
  k_{\max}(p)\;=\;\max\bigl\{k\ge0:\ (p,k)\ \text{is marked}\bigr\},
\]
with the convention that $(p,0)$ is always marked, so that the marks are
nested by construction and no Poisson law is available for the full set
of marked pairs.  Write
\[
  \kappa(r,k;p)\;=\;\lim_{H\to\infty}\mathbb
  E\bigl[p^{rk}/\Delta_{\Gamma(H)}(p^{k})\bigr]
\]
for the limiting expected local Kummer factor of the sampled tuple at
depth $k$, computed from the Kummer degrees of
Section~\ref{sec:notation} and equal to $1$ when almost every sampled
tuple is Kummer-generic at $p$.  Existence of the limit, and hence
independence of the intensity from the height parameter $H$, is part of
the assertion.  Then:
\emph{(i)} after the explicit Kummer local factors are inserted, the
process of maximal-depth points
$\bigl\{\bigl(p,k_{\max}(p)\bigr):k_{\max}(p)\ge1\bigr\}$ converges as
$H\to\infty$ to a Poisson process on pairs $(p,k)$ with intensity
\[
  \kappa(r,k;p)\,p^{-rk}\;-\;\kappa(r,k+1;p)\,p^{-r(k+1)} ,
\]
equivalently, for each fixed $k\ge1$ the depth-$k$ slice $\{p:
k_{\max}(p)\ge k\}$ converges to a Poisson process of intensity
$\kappa(r,k;p)p^{-rk}$.  The Poisson statement is made for the
maximal-depth process and for the individual depth slices, and not for
the nested family of marked pairs.
\emph{(ii)} conditionally on a depth-$k$ mark at $p$, the depth-$(k+1)$
mark is a Haar lift imposing $r$ further linear conditions, so the depths
form a codimension-$r$ tower rather than an independent sequence, which
is the deterministic nesting recorded above.
\emph{(iii)} for a fixed $\Gamma$ the same intensities give the
Borel--Cantelli predictions of Conjecture~\ref{c13}.  Clause (iii) is a
conjecture and not a consequence of (i), since a single fixed $\Gamma$ is
not a sample from the height ensemble and carries no probability space of
its own.
\end{conjecture}

The $p$-adic logarithm of $r$ independent generators is a vector in a
rank-$r$ lattice and each further digit imposes another codimension-$r$
condition, while the height ensemble supplies the probability space that
a fixed group cannot.  The statement predicts all higher-depth
exceptional strata, the dependence between depths, and the law of the
finite exceptional set when $rk>1$, and it connects Wieferich primes to
the theory of random $p$-adic Smith normal forms.  Higher Wieferich
congruences are usually treated one base and one depth at a time, in the
regulator language of Boeckle, Guiraud, Kalyanswamy and Khare or the
$abc$-conditional language of Silverman and of Graves and Murty, and the joint
rank-by-depth point process with a declared family ensemble is what is
claimed here.

\emph{Evidence.}  No base-$2$ prime of depth at least two occurs below
$250{,}000$ while the depth-one primes are $1093$ and $3511$, and since
$\sum_{p\ge3}p^{-2}\approx0.202$ the stated intensity predicts about
$0.2$ such primes in that range, so the observation is compatible with
clause (i) and tests nothing about the family ensemble on which the
conjecture is built.

\begin{conjecture}[Toric finite logarithm and codimension
principle]\label{c15}
Let $T/\mathbb Q$ be an algebraic torus with connected N\'eron model over
$\mathbb Z[1/N]$, let $P\in T(\mathbb Q)$, and let $S$ be the identity
component of the Zariski closure of the subgroup generated by $P$, of
dimension $r$.  Write $\mathcal T$ for the connected N\'eron model.  A
prime $p$ is \emph{good} if $p\nmid N$, if $p$ is unramified in the
splitting field of $T$, and if $P$ is integral at $p$, that is if
$P\in T\bigl(\mathbb Z[1/N]\bigr)$ so that $p$ divides no denominator of
$P$ and $P$ extends to a section of $\mathcal T$ over $\mathbb
Z_{(p)}$.  For a good prime $p$ let $N_p$ be the exponent of $T(\mathbb
F_p)$.  Then $[N_p]P$ reduces to the identity modulo $p$, so its image in
$\mathcal T(\mathbb Z/p^{2}\mathbb Z)$ lies in the kernel of reduction,
and the first finite logarithm $\ell_p(P)$ is defined as its image under
the canonical isomorphism
\[
  \ker\bigl(\mathcal T(\mathbb Z/p^{2}\mathbb Z)\longrightarrow
  \mathcal T(\mathbb F_p)\bigr)\ \xrightarrow{\ \sim\ }\
  \Lie(\mathcal T)(\mathbb F_p)\;=\;\Lie(T)(\mathbb F_p),
\]
the isomorphism of the first congruence neighbourhood of the identity
with the Lie algebra, which for $T=\mathbb G_m$ sends $1+pu$ to $u$ and
recovers the Fermat quotient.  Then:
\emph{(i)} [elementary] $\ell_p(P)\in\Lie(S)(\mathbb F_p)$, so the
algebraic relations satisfied by $P$ annihilate the corresponding
tangent coordinates identically.
\emph{(ii)} conditional on any compatible Frobenius class in the
splitting field of $T$, the normalized vector $\ell_p(P)/p$, read in the
reduction of a fixed $\mathbb Z$-basis of $\Lie(S)$, equidistributes with
respect to Haar measure on $(\mathbb R/\mathbb Z)^{r}$.
\emph{(iii)} [functoriality] for a morphism of tori $\varphi\colon T\to
T'$ one has $\ell_p(\varphi(P))=d\varphi\bigl(\ell_p(P)\bigr)$ at every
good prime for both, so the construction and the law are functorial and
the codimension invariant is geometric rather than a coordinate count.
\emph{(iv)} vanishing of $\ell_p(P)$ to $p$-adic depth $k$ has local
density $p^{-rk}$ times the explicit splitting factor determined by the
decomposition of $p$ in the splitting field of $T$, the density being
that of the target inside $\Lie(S)(\mathbb Z/p^{k}\mathbb Z)$ under Haar
measure.  This is a statement about the measure of the target and not a
probability over a declared ensemble of primes, no such ensemble being
attached to a fixed $P$, and the divergent-summable boundary of the
associated Borel--Cantelli sum $\sum_pp^{-rk}$ therefore depends only on
$\dim S$ and $k$.
\emph{(v)} the case $T=\mathbb G_m$, and more generally the case of
powers of $\mathbb G_m$, recovers Conjectures~\ref{c11} and~\ref{c13}.
The construction depends on the integral model, and the statement is
asserted for the connected N\'eron model at good primes, where the
component group is trivial.
\end{conjecture}

The ordinary Fermat quotient is the case of the multiplicative group,
and replacing coordinate counting by the tangent space of the algebraic
subgroup generated by the point turns the exceptional exponent into a
geometric invariant that transforms correctly under morphisms.  A proof
would extend Wieferich and mod-$p$ regulator heuristics to norm-one tori,
to units in number fields and to higher-dimensional algebraic groups, and
it predicts which relations survive base change and which exceptional
prime problems are genuinely of higher rank.  The finite-logarithm
picture for the multiplicative group is due in its analytic form to
Shparlinski and in its regulator form to Boeckle, Guiraud, Kalyanswamy and
Khare, and the extension to arbitrary tori with dimension and
functoriality controlling the exponent is the broader framework claimed
here.

\emph{Evidence.}  For the norm-one torus of $\mathbb Q(\sqrt2)$ with
$P=3+2\sqrt2$ and $22{,}041$ primes, the scalar tangent component
vanished identically as the norm relation of clause (i) requires, the
fixed Fourier magnitudes of the one-dimensional tangent coordinate lay
between $0.0034$ and $0.0099$, the split and inert prime means were
$0.4975$ and $0.5000$, and exact toric vanishing occurred at $p=13$ and
$p=31$, all of which are fixed-frequency consistency checks on clause
(ii) for a single point on a single torus.

\begin{question}\label{q:toric}
For a norm-one torus of a real quadratic field, is the local splitting
factor of Conjecture~\ref{c15}(iv) at inert primes, where $|T(\mathbb
F_p)|=p+1$, equal to the split-prime factor, or does the difference of
group orders shift it?  This is an open question and is not counted among
the twenty-five conjectures.
\end{question}

\section{Programme IV. Arboreal resolution and primitive dynamical
factors}\label{sec:IV}

An entropy-conductor boundary controls how much of an infinite arboreal
representation is visible below a given height, and the primitive factors
of orbit terms and the gcds between orbits are governed by arboreal and
geometric structure.  The family-averaged clauses of this programme use
the parameter law (P) of Section~\ref{sec:notation}.

\begin{conjecture}[Joint arboreal Chebotarev law at growing
depth]\label{c16}
Let $f\in\mathbb Q[x]$ be a non-postcritically-finite polynomial of
degree $d\ge2$ and let $a\in\mathbb Q$ have infinite backward orbit and
arboreal image $G_\infty$ of finite index in $\Aut(T_\infty)$.  Let
$n=n(X)$ satisfy $E(n,X)\to0$ in the notation of
Conjecture~\ref{c17}, and sample $p$ uniformly from the primes $p\le X$
unramified in $K_n$.  Then:
\emph{(i)} the complete compatible vector of factorization types of
$f^{j}(x)-a$ modulo $p$ for every $j\le n$ converges in total variation
to the Haar--Chebotarev distribution on conjugacy-cylinder data of $G_n$.
\emph{(ii)} every statistic obtained by projecting the tree, in
particular root counts at each level, cycle lengths, first-preimage
levels and cross-level coincidence events, converges jointly and
uniformly in the same regime, so the convergence is at the level of the
whole visible tree and not one quotient at a time.
\end{conjecture}

Chebotarev supplies one fixed quotient at a time, and the content here is
that the entire visible finite tree stays equidistributed until the
combined sample and conductor complexity reaches the critical value
identified in Conjecture~\ref{c17}.  A proof would turn an infinite
profinite representation into a quantitatively observable arithmetic
process, transferring group martingales and cycle statistics to growing
finite-field depth, and it is the foundation on which
Conjectures~\ref{c18} to~\ref{c20} rest.  Arboreal representations
control fixed-level factorization and prime-divisor densities in the work
surveyed by Jones and in the results of Doyle, Healey, Hindes and Jones,
with fixed-point proportions computed by Juul, Kurlberg, Madhu and Tucker
and by Fari\~na-Asategui and Radi, and the joint law for all levels
simultaneously in the maximal resolvable range is the claimed content.

\emph{Evidence.}  For quadratic maps $f_c(x)=x^2+c$ with $c$ drawn from a
list of eleven parameters and primes $7\le p\le499$, the root-count
distributions at levels $2$, $3$ and $4$ lay within total variation
$0.042$ to $0.048$ of the full binary wreath-product model, but the
census covers ten effective maps rather than eleven, because the
good-prime filter used, $p\not\equiv0\pmod c$, removes the parameter
$c=1$ entirely and otherwise deletes only the single prime $p=7$ at
$c=\pm7$, and the test is fixed-level and therefore subcritical.

\begin{conjecture}[Sharp entropy-conductor barrier for growing-level
Chebotarev]\label{c17}
Let $K_n$, $G_n$, $D_n$ and $E(n,X)$ be as in
Section~\ref{sec:notation}, and let $\mathrm{TV}(n,X)$ denote the total
variation distance between the empirical frequency vector of
conjugacy cylinders up to level $n$, computed over the primes $p\le X$
unramified in $K_n$, and the Haar--Chebotarev vector.  Write
\[
  D_n(f,X)\;=\;\max_{C}\left|
  \frac{\#\{p\le X\ \text{unramified in}\ K_n:\ \mathrm{Frob}_p\in
  C\}}{\#\{p\le X\ \text{unramified in}\ K_n\}}
  \;-\;\frac{|C|}{|G_n|}\right|
\]
for the \emph{maximal cylinder discrepancy} at level $n$, the maximum
being over the conjugacy cylinders $C$ of $G_n$.  Call
$(f_c)_{|c|\le H}$ a \emph{one-parameter family of finite-index type} if
every $f_c$ has degree $d\ge2$, if the basepoint has infinite backward
orbit for every $c$ outside a set of density zero in the parameter law
(P), and if for every such $c$ the arboreal image
$G_\infty(f_c)$ has finite index in $\Aut(T_\infty)$ with the index
bounded uniformly in $c$.  Then:
\emph{(i)} [upper half, conditional proposition under the generalized
Riemann hypothesis for the Dedekind zeta functions of the fields $K_n$]
one has $\mathrm{TV}(n,X)\ll E(n,X)^{1/2}(\log X)^{O(1)}$, the first term
of $E$ supplying the multinomial sampling error over $|G_n|$ states and
the second the effective Chebotarev error summed over states, so
$E(n,X)(\log X)^{O(1)}\to0$ implies $\mathrm{TV}(n,X)\to0$.  The
logarithmic factor is not decorative.  The effective Chebotarev bound
under a Riemann hypothesis carries powers of $\log D_n$ and of $\log X$,
and the second term of $E$ records the square-root conductor obstruction
without them, so clause (i) is asserted only up to factors of that
shape.  The sharpness pairing with clause (ii) is unaffected, since the
two clauses are separated by a power of $E$ and not by a logarithm.
\emph{(ii)} [lower half, family-averaged] let $(f_c)_{|c|\le H}$ be a
one-parameter family of finite-index type in the sense above, with $c$
sampled by the parameter law (P), and suppose $E(n,X)\to\infty$ along a
sequence of pairs $(n,X)$.  Then the family average $\mathbb
E_c\bigl[D_n(f_c,X)\bigr]$ of the maximal cylinder discrepancy at level
$n$ stays bounded below by a positive absolute
constant along that sequence.  The clause is stated as a family average
because a single fixed $f$ produces a deterministic sequence of
discrepancies for which no probability space is declared, and an
``infinitely often'' assertion for one $f$ is therefore not the form
asserted here.
\emph{(iii)} [critical window, conjectural profile] in the window
$E(n,X)\asymp1$ the centred and normalized empirical cylinder-frequency
field has a Gaussian-multinomial limit whose covariance is the
multinomial covariance projected by the tree quotient relations.  This
clause is the most speculative of the three and no finite computation
below bears on it.
The sharpness assertion of this conjecture is the pairing of (i) with
(ii): the analytic input makes the subcritical side a deduction, and the
conjectural content is that the same functional governs the supercritical
side on average over a family.
\end{conjecture}

The first term of $E(n,X)$ is the information-theoretic requirement that
every state of $G_n$ receive enough primes, the second is the square-root
Chebotarev and conductor obstruction, and the claim is that their maximum
rather than the tree level alone is the canonical resolution parameter.
A sharp barrier would say exactly how much of an infinite arboreal
representation can be observed below a given height, would organize
growing-level questions throughout arithmetic dynamics, and would
separate group growth from ramification growth as sources of
irresolvability.  The effective Chebotarev bounds that make clause (i) a
deduction are those of Thorner and Zaman in the tradition of Lagarias and
Odlyzko, while arboreal group calculations in the surveys of Jones and in
Doyle, Healey, Hindes and Jones fix the growth of $|G_n|$, and the
identification of the maximal resolvable depth is the claimed content.

\emph{Evidence.}  Only levels $2$ to $4$ were computationally accessible,
over ten effective quadratic maps rather than the eleven parameters
listed, since the good-prime filter removes $c=1$ entirely, and their
root-count distributions lay within total variation $0.042$ to $0.048$ of
the wreath-product law, which is subcritical consistency alone and bears
on neither the sharpness pairing of clauses (i) and (ii) nor the critical
profile of clause (iii).

\begin{conjecture}[Poisson--Dirichlet law for primitive factors in
dynamical families]\label{c18}
Fix $d\ge2$ and sample $c$ by the parameter law (P), uniformly from the
integers with $0<|c|\le H$ outside the postcritically finite and
reducible loci.  Let $n=n(H)$ grow so slowly that $d^{\,n}=o(\log H)$,
and let $P_n(c)$ be the primitive part of $v_n(c)=f_c^{n}(0)$ in the
sense of Section~\ref{sec:notation}, that is, the largest divisor
composed of primes dividing no earlier orbit term.  Then:
\emph{(i)} the ranked ratios $\log\ell_i/\log P_n(c)$, over the primes
$\ell_i$ dividing $P_n(c)$ counted with valuation, converge in
distribution as $H\to\infty$ to $\PD(1)$.
\emph{(ii)} the limit is jointly marked by the level of first appearance
of each prime and by its Frobenius data in the arboreal tower of
Conjecture~\ref{c16}, the marks being conditionally independent given the
masses.
The growth condition $d^{\,n}=o(\log H)$ and the family sampling are part
of the statement, since a fixed parameter would give a deterministic
sequence of factorizations with no probability space attached.
\end{conjecture}

Once the inherited factors forced by rigid divisibility are deleted, the
primitive part should behave as a well-distributed large integer with no
further macroscopic algebraic decomposition, and $\PD(1)$ is the
canonical scale-invariant factor process for such an integer.  A proof
would connect arithmetic dynamics to modern factor-process theory and
would quantify far more than the existence of a primitive divisor,
predicting largest factors, smoothness and the number of macroscopic new
primes.  Primitive divisors and average Zsigmondy sets are established
themes, in Zsigmondy's classical theorem and in the work of Ingram and
Silverman and of Laishram, Rout and Yadav, and large prime factors of
well-distributed sequences carry Poisson--Dirichlet behaviour by
Bharadwaj and Rodgers, so the claimed content is the $\PD(1)$
factorization of the primitive part in a dynamical family with arboreal
marks.

\emph{Evidence.}  For $158$ parameters at level $4$ the mean largest
logarithmic share was $0.684$ with median $0.658$ against the
Golomb--Dickman mean $0.6243$ for $\PD(1)$, a gap of several standard
errors of the mean if the parameters are treated as independent samples,
so the finite-level structure is visible rather than absent, and level
$4$ with $d^{\,n}=16$ against $\log H\approx4.4$ is far outside the
stated growth regime.

\begin{conjecture}[Euler law for squarefull primitive orbit
factors]\label{c19}
In the family and growth regime of Conjecture~\ref{c18}, let
$\delta_p(n)$ be the density of residues $c\bmod p^{2}$ for which
$p^{2}\mid v_n(c)$ while $p$ divides no earlier orbit term, a quantity
computed by counting parameters modulo $p^{2}$ and therefore derived
rather than fitted.  Let the \emph{critical-orbit collision locus} at
level $n$ be the set of parameters
\[
  \mathcal R_n\;=\;\bigl\{c:\ v_i(c)\equiv v_j(c)\ (\mathrm{mod}\ p^{2})
  \ \text{for some}\ 1\le i<j\le n\ \text{and some}\ p\mid P_n(c)\bigr\},
\]
on which two distinct forward images of the critical point $0$ agree to
the second power of a primitive prime and the codimension-two congruence
defining a repeated primitive factor degenerates.  Then:
\emph{(i)} the sequence $\delta_p(n)$ stabilizes, in the sense that there
is $n_0(p)$ with $\delta_p(n)=\delta_p$ for every $n\ge n_0(p)$, and with
that stable value the probability that $P_n(c)$ is squarefree converges
to $\prod_p(1-\delta_p)$.
\emph{(ii)} after conditioning on the complement of $\mathcal R_n$, the
valuations $v_p(P_n(c))$ at distinct primes are asymptotically
independent.
\emph{(iii)} the logarithm of the squarefull part of $P_n(c)$ is tight as
$H\to\infty$.
\end{conjecture}

A repeated primitive factor is a codimension-two congruence condition on
the parameter unless ramification in the critical orbit lowers the
codimension, so the exceptional ramified loci must be separated
explicitly before the Euler product can be asserted.  The statement
supplies the $p$-adic input that makes Conjecture~\ref{c18} sound, since
a systematic square factor would distort the mass normalization, and it
refines dynamical Zsigmondy theory from an existence statement to a
multiplicity statistic.  Squarefree-value heuristics are classical, and
the primitive and inherited decomposition together with the
critical-orbit correction is what makes the law specifically dynamical,
in the setting of Ingram and Silverman and of Laishram, Rout and Yadav.

\emph{Evidence.}  At level $4$, $155$ of the $158$ tested primitive parts
were squarefree, a fraction $0.981$, and the three repetitions are
recorded rather than discarded, but $158$ samples at a single level
cannot distinguish the conjectured Euler product from nearby constants
and give no information about clauses (ii) and (iii).

\begin{conjecture}[Geometric-arboreal criterion for large dynamical
gcds]\label{c20}
Let $f,g\in\overline{\mathbb Q}[x]$ be disintegrated of degree at least
two and let $a,b$ have positive canonical height for $f$ and $g$
respectively.  For $\alpha,\beta\in\overline{\mathbb Q}^{\times}$ lying
in a number field $K$, the \emph{logarithmic gcd} is Silverman's
generalized logarithmic greatest common divisor
\[
  \log\gcd(\alpha,\beta)\;=\;\sum_{v\in M_K^{0}}n_v\,
  \min\bigl\{v^{+}(\alpha),\,v^{+}(\beta)\bigr\},
  \qquad n_v=\frac{[K_v:\mathbb Q_v]}{[K:\mathbb Q]},
\]
the sum being over the finite places of $K$ with
$v^{+}(x)=\max\{0,-\log|x|_v\}$, a quantity independent of the field $K$
containing $\alpha$ and $\beta$ and equal to $\log\gcd(a,b)$ for positive
integers.  Consider
\[
  \limsup_{\min(n,m)\to\infty}\
  \frac{\log\gcd\bigl(f^{n}(a),g^{m}(b)\bigr)}
  {\max\bigl(h(f^{n}(a)),\,h(g^{m}(b))\bigr)}\;>\;0 ,
\]
the limit superior being taken over all pairs $(n,m)$ of positive
integers with $\min(n,m)\to\infty$, so that both orbits advance.
Then:
\emph{(i)} [constructive direction] if, after replacing $f$ and $g$ by
iterates, the orbit pair lies on a nontrivial periodic algebraic
correspondence for $(f,g)$, equivalently if $f$ and $g$ have a common
semiconjugate dynamical quotient carrying the branches of $a$ and $b$,
then the displayed limsup is positive.
\emph{(ii)} [strong converse layer] conversely, positivity of the limsup
forces the existence of such a correspondence, so no accidental
congruence source of macroscopic gcds exists.  Clause (ii) is separately
falsifiable from (i), and a thin infinite family of large gcds produced
by a correspondence visible only after base change, or by a nonperiodic
special subvariety, refutes it alone.
\emph{(iii)} [functoriality] such a correspondence induces a common
quotient of the distinguished arboreal representations of $f$ and $g$.
\emph{(iv)} in the absence of a correspondence the normalized gcd is
$o(1)$ outside a set of pairs $(n,m)$ of density zero, and the
deterministic second-moment bound
\[
  \frac1{N^{2}}\sum_{n,m\le N}
  \left(\frac{\log\gcd\bigl(f^{n}(a),g^{m}(b)\bigr)}
  {\max\bigl(h(f^{n}(a)),\,h(g^{m}(b))\bigr)}\right)^{2}
  \;\ll\;\frac1{\log N}
\]
holds as $N\to\infty$.  The bound is deterministic and is asserted over
the pairs $(n,m)$ directly.  The pairs carry no sampling law, so no
covariance, connected or otherwise, is attributed to the sequence of
common primes.
\end{conjecture}

A common prime is a Frobenius event in the fibre product of the two
preimage towers, but only the common quotients that arise from geometry
persist strongly enough to create macroscopic gcds, which is why a
criterion phrased through finite quotients alone is too broad.  An
equivalence would connect dynamical unlikely intersections, gcd bounds
and profinite Galois representations in one statement, and either outcome
identifies the responsible invariant, geometry rather than accidental
congruence.  Orbit gcds, invariant curves and arboreal entanglement are
studied separately in the work of Ingram and Silverman, of Asgarli and
coauthors on collisions of orbits over function fields, and in the
classification of invariant varieties by Medvedev and Scanlon, which is
what makes the disintegrated hypothesis the right one.  The
gcd bounds for values of algebraic functions are those of Bugeaud,
Corvaja and Zannier, and the dynamical gcd conjectures of Silverman are
the immediate predecessors of the constructive direction, so the claimed
content is that the geometric correspondence is the complete source.

\emph{Evidence.}  For small quadratic maps at levels $2$ to $6$ the
independent-parameter pairs had mean maximum normalized logarithmic gcd
$0.035$ with maximum $0.069$ whereas deliberately related same-orbit
pairs averaged $0.097$ and reached $0.246$, which illustrates the
predicted separation over nine parameters and five levels and tests
neither direction of the equivalence.

\section{Programme V. Frobenius-marked polynomial
factorization}\label{sec:V}

Macroscopic factors of polynomial values carry Galois marks, the small
and large scales decouple once the correct normalization is used,
reducible polynomials glue their colour processes through an explicit
local factor, and smoothness and primality are joined by a sieve flow.
Every statement uses the uniform law on $n\le N$ declared in
Section~\ref{sec:notation}.

\begin{conjecture}[Frobenius-marked Poisson--Dirichlet factor
process]\label{c21}
Let $f\in\mathbb Z[x]$ be irreducible with splitting field $L$ and
splitting group $G=\Gal(L/\mathbb Q)$ acting transitively on the roots.
Sample $n$ uniformly from $\{1,\dots,N\}$, mark each prime factor $p$ of
$f(n)$ by its Frobenius conjugacy class $C$ in $G$, and give it mass
$u=\log p/\log|f(n)|$.  Then:
\emph{(i)} the macroscopic marked process converges as $N\to\infty$ to a
marked $\PD(1)$ process whose mark intensity is proportional to
$(|C|/|G|)\fix(C)\,du/u$, the normalization being exactly one by
Burnside's lemma, so no constant is fitted.
\emph{(ii)} for every fixed number of macroscopic factors, their
Frobenius marks are asymptotically independent conditional on their
masses, once the finite local congruence state of $n$ is removed.
\emph{(iii)} [functoriality] the law is functorial under quotienting the
splitting group, the marked process for a quotient $G/H$ being the
pushforward of the marked process for $G$.
\end{conjecture}

A prime divides a value of $f$ with frequency equal to its number of
roots modulo $p$, so Chebotarev size-biases the conjugacy classes by
fixed-point counts while scale invariance of the factorization produces
the Poisson--Dirichlet masses, and the two mechanisms compose without a
free parameter.  A proof would unify large-factor statistics with Galois
theory, predicting not only how large the macroscopic factors are but how
they split in the defining field, and it extends naturally to covers,
norm forms and arithmetic monodromy families.  Large prime factors of
sufficiently well-distributed sequences have Poisson--Dirichlet behaviour
by Bharadwaj and Rodgers, with the prime-value side developed by Kravitz,
Woo and Xu, and the simultaneous Galois mark on every macroscopic factor
is the claimed content.

\emph{Evidence.}  For $f(x)=x^3-2$ and $n\le1200$ the largest prime
factor had one root modulo $p$ in $50.58$ per cent of samples and three
roots in $49.42$ per cent against the size-biased prediction of one half
and one half for $S_3$, a deviation of $0.58$ percentage points against a
descriptive standard error of about $1.4$ percentage points at that
sample size, so the comparison is a consistency check at a single degree
of freedom and bears on neither clause (ii) nor clause (iii).

\begin{conjecture}[Small-large factor independence after logarithmic mass
removal]\label{c22}
Let $f\in\mathbb Z[x]$ be irreducible, sample $n$ uniformly from
$\{1,\dots,N\}$, and for $y=N^{o(1)}$ let $S_y(f(n))$ be the $y$-smooth
part and $R_y=f(n)/S_y(f(n))$ the residual cofactor.  Then:
\emph{(i)} the small-prime Kubilius field generated by the valuations
$v_p(f(n))$ for $p\le y$ is asymptotically independent of the
Frobenius-marked $\PD(1)$ process of $R_y$, provided the latter is
normalized by $\log R_y$ and not by $\log|f(n)|$.
\emph{(ii)} [strong converse layer] the only leading coupling between the
two scales is the removed logarithmic mass together with the explicit
resultant data of $f$, so no further interaction survives.  Clause (ii)
is separately falsifiable from (i): residual dependence transmitted
through Frobenius, or through values unusually close to smooth numbers,
refutes it alone.
\end{conjecture}

Small-prime divisibility consumes logarithmic mass and therefore
mechanically depresses the raw largest-factor share, so unqualified
independence of the small and large scales is false by mass conservation,
and the substantive statement is that once the remaining mass is used as
the denominator, scale invariance restores independence.  A proof would
be a multiscale factorization principle allowing local sieve information
and global factor-shape information to be combined without double
counting, which is exactly the step that heuristics for polynomial values
usually take without justification.  The Poisson--Dirichlet side is the
theory of Bharadwaj and Rodgers, and the corrected conditional
independence with the residual-mass normalization is the claimed content.

\emph{Evidence.}  For $n^3-2$ the raw Pearson correlation between the
number of prime factors at most $100$, counted with multiplicity, and the
largest-factor logarithmic
share was $-0.465$ and fell to $0.104$, with Spearman $0.074$, after
renormalizing by the logarithm of the residual cofactor, so the
correction is empirically essential, though $0.104$ sits about three
descriptive standard errors from zero at roughly $1200$ samples and is
therefore consistent with a slowly vanishing finite-$N$ coupling rather
than with exact independence at this range.

\begin{conjecture}[Coloured factor superposition with local-exclusion
gluing]\label{c23}
Let $f=\prod_{i=1}^{s}g_i\in\mathbb Z[x]$ be squarefree with irreducible
components $g_i$, sample $n$ uniformly from $\{1,\dots,N\}$, and let
$y=N^{o(1)}$.  Let $\rho(I,p)=\#\{x\bmod p:\ p\mid g_i(x)\ \text{for all}
\ i\in I\}$ be the joint root-count tensor.  Then:
\emph{(i)} conditioned on the full joint local state
$\bigl(v_p(g_i(n))\bigr)_{i\le s,\ p\le y}$, and with each component
normalized by its own residual logarithmic mass, the coloured
Frobenius-marked macroscopic factor processes of the $g_i$ are
asymptotically independent $\PD(1)$ processes.
\emph{(ii)} before conditioning, their entire interaction is the explicit
Gibbs--Euler factor assembled from the tensors $\rho(I,p)$, whose leading
part is
\[
  \mathcal G(f)\;=\;\prod_p\frac{\displaystyle\sum_{I\subseteq\{1,\dots,s\}}
  (-1)^{|I|}\frac{\rho(I,p)}{p}}
  {\displaystyle\prod_{i=1}^{s}\Bigl(1-\frac{\rho(\{i\},p)}{p}\Bigr)},
\]
the numerator being the density of residues modulo $p$ avoiding every
component by inclusion and exclusion and the denominator the same density
computed as if the components were independent, so that the primes
dividing the resultant contribute hard shared atoms and the remaining
primes contribute exclusion or attraction according to whether
$\rho(I,p)$ falls below or above the independent product.  [stated
programme] The closed assembly of the full Gibbs weight from the tensors
$\rho(I,p)$ at every $|I|$, of which $\mathcal G(f)$ is the two-point
part, is the derivation route registered with this conjecture.
\emph{(iii)} [strong converse layer] no further long-range colour
coupling survives.  Clause (iii) is separately falsifiable from (i) and
(ii), and an archimedean size correlation or a medium-prime coupling
persisting after every fixed-$y$ conditioning refutes it alone.
\end{conjecture}

Algebraic components generate separate macroscopic masses, but the same
residue $n$ modulo $p$ couples their small-prime divisibility even when
the resultant is one, so a reducible polynomial is not a product of
independent random integers and deleting only the resultant primes is not
enough.  The statement supplies the correct functorial rule under
multiplication of polynomials, predicting coloured largest factors, joint
smoothness and common-cofactor risk without the false independence
assumption that simpler heuristics use.  The Poisson--Dirichlet
background is again that of Bharadwaj and Rodgers, and the identification
of the complete interaction as a Gibbs--Euler factor built from joint
root-count tensors is what is claimed.

\emph{Evidence.}  For $g_1=n^2+1$ and $g_2=n^2+n+1$, whose resultant is
one, no common prime factor occurred in $6000$ samples and the
largest-factor logarithmic shares had correlation $-0.027$, about two
descriptive standard errors from zero, but the computation does not
condition on the joint local state, so clauses (i) and (iii) are untested
and only the macroscopic null case predicted by clause (ii) at resultant
one is observed.

\begin{conjecture}[Saddle-point entanglement law for joint polynomial
smoothness]\label{c24}
Let $f_1,\dots,f_s\in\mathbb Z[x]$ be pairwise coprime and primitive,
sample $n$ uniformly from $\{1,\dots,N\}$, and let the smoothness
parameters $u_i=\log|f_i(n)|/\log y_i$ range over the \emph{admissible
region}, by which is meant a fixed compact subset of
$\{(u_1,\dots,u_s):u_i>1\ \text{for every}\ i\}$ together with the
restriction $y_i\ge(\log N)^{1+\varepsilon}$ for a fixed
$\varepsilon>0$, the range in which the saddle-point expansion for smooth
values is available.  For $J\subseteq\{1,\dots,s\}$ and a prime $p$ let
\[
  \mu_p^{\sigma}(J)\;=\;\frac{\rho(J,p)\,p^{-\sum_{i\in J}\sigma_i}}
  {\sum_{J'\subseteq\{1,\dots,s\}}\rho(J',p)\,p^{-\sum_{i\in J'}\sigma_i}}
\]
be the tilted joint local measure attached to a vector
$\sigma=(\sigma_1,\dots,\sigma_s)$, with $\rho(J,p)$ the joint root-count
tensor of Section~\ref{sec:notation}.  The \emph{joint saddle equations}
are the $s$ coupled equations
\[
  \sum_{p\le y_i}\ \sum_{J\ni i}\mu_p^{\sigma}(J)\,\log p
  \;=\;\log|f_i(n)|,\qquad 1\le i\le s,
\]
which determine $\sigma$ and reduce for $s=1$ to the classical
saddle-point equation for smooth numbers.  Then:
\emph{(i)} the joint density of $n$ for which every $f_i(n)$ is
$y_i$-smooth equals the product of the component Dickman probabilities
multiplied by an Euler product of the tilted local root distributions
$\mu_p^{\sigma}$ at the solution $\sigma$ of the joint saddle equations,
so the entanglement factor is a function of $(u_1,\dots,u_s)$ and not a
fixed constant.
\emph{(ii)} [strong converse layer] after this tilt no additional
long-range correlation between the components survives.  Clause (ii) is
separately falsifiable from (i), and a correlation persisting near the
edge of the smoothness range, or an archimedean factor not captured by
the Euler product, refutes it alone.
\end{conjecture}

Conditioning on smoothness changes the effective probability that each
small prime divides a value, so the local dependence between components
must be computed under the tilted measure rather than under the
unconditioned Euler product, and that is precisely why a fixed correction
factor multiplying Dickman probabilities is too crude.  A proof would
give a coherent local-to-global theory of simultaneous smooth polynomial
values and would correct a common heuristic error, linking saddle-point
analysis, resultants and multivariate sieves.  The single-polynomial
theory rests on Dickman and on the saddle-point method for smooth
numbers, in the uniform form developed by Hildebrand and Tenenbaum, and
the parameter-dependent entanglement factor is the claimed content.

\emph{Evidence.}  For $n^2+1$ and $n^2+n+1$ the ratio of the joint
probability to the product of the marginals was $0.674$, $0.893$ and
$0.952$ as the largest-factor threshold rose from $0.35$ to $0.55$, and
the monotone movement of that ratio is what rejects a single fixed
correction at this range, though the thresholds are largest-factor shares
rather than genuine smoothness parameters and the sample is $6000$
values.

\begin{conjecture}[Buchstab--Bateman--Horn sieve-to-prime
bridge]\label{c25}
Let $f\in\mathbb Z[x]$ be primitive, irreducible and admissible, sample
$n$ uniformly from $\{1,\dots,N\}$, and put $u=\log|f(n)|/\log y$ and
$V_f(y)=\prod_{p\le y}\bigl(1-\rho_f(p)/p\bigr)$.  Then:
\emph{(i)} uniformly for $u$ in compact subsets of $(1,\infty)$, the
density of $y$-rough values, those with no prime factor at most $y$, is
$V_f(y)\,\e^{\gamma}\omega(u)\bigl(1+o(1)\bigr)$.
\emph{(ii)} conditional on $y$-roughness, the probability that $f(n)$ is
prime is
\[
  \frac{C_f/\log|f(n)|}{V_f(y)\,\e^{\gamma}\omega(u)}\bigl(1+o(1)\bigr),
\]
with $C_f$ the Bateman--Horn constant of $f$.
\emph{(iii)} [elementary endpoint] for $1<u\le2$ a $y$-rough value below
$y^2$ has no admissible nontrivial factorization and is therefore prime
unless it is a unit, so the conditional probability in (ii) equals one
exactly on that range, the values being positive on the range considered
and the unit case therefore not arising.
The conditional probabilities of (ii) are consistent by construction as
$y$ increases, each being the ratio of the same Bateman--Horn density to
the roughness density at that scale, and no martingale property is
asserted.
\end{conjecture}

The finite sieve accounts for the local factors up to $y$, the Buchstab
evolution accounts for the number of remaining factors, and Bateman--Horn
supplies the terminal one-factor density, so the parity barrier appears
as a boundary condition at $u\le2$ rather than as an unexplained
constant.  A proof would reframe Bateman--Horn as a scale-by-scale law
and would provide a concrete interface between sieve theory and
prime-producing polynomials, predicting rough values, semiprime
transitions and primality in a single equation.  Bateman and Horn predict
the terminal density, Buchstab theory predicts the finite-scale
roughness, and the uniform conditional bridge between them, including the
parity endpoint, is the distinctive statement.

\emph{Evidence.}  For $f(n)=n^2+1$ on $60{,}000\le n\le140{,}000$ and
$y=13,31,71,151,313$ the observed rough-survivor densities agreed with
the finite-$u$ prediction of clause (i) within $0.11$ per cent and the
observed conditional prime probabilities were uniformly $0.984$ to
$0.985$ of the prediction of clause (ii) with no parameter fitted, and
that uniform shortfall of about $1.5$ per cent across five values of $y$
is a systematic secondary term rather than sampling noise.  The five
rows are not five independent measurements: they share one nested sample
of $n$, the rough survivors at a larger $y$ being a subset of those at a
smaller one, so the five agreements are strongly correlated and the
consistency of the shortfall across the rows carries much less
information than five independent confirmations would.

\section{Reading the collection}\label{sec:reading}

The twenty-five statements are grouped by mechanism rather than by
difficulty.  Conjectures~\ref{c1}, \ref{c3}, \ref{c6}, \ref{c7},
\ref{c11}, \ref{c16}, \ref{c18}, \ref{c21}, \ref{c24} and~\ref{c25} are
parent mechanisms, each supplying the objects that its programme then
uses.  Conjectures~\ref{c2}, \ref{c4}, \ref{c14}, \ref{c19}
and~\ref{c22} are independent structural extensions that introduce a
cluster invariant, a valuation process or a conditional law not contained
in the parent.  Conjectures~\ref{c8}, \ref{c13} and~\ref{c17} are phase
boundaries.  Conjecture~\ref{c9} is a sharp anomaly statement, and
Conjectures~\ref{c5}, \ref{c10}, \ref{c12}, \ref{c15}, \ref{c20}
and~\ref{c23} are converse classifications or functorial laws.

Eleven conjectures carry a strong converse layer that is stated
separately from its constructive companion and can fail on its own.  The
eleven tags are of one kind and form one list: Conjecture~\ref{c1}(ii) on
the absence of coefficients outside the overlap-connected expansion,
Conjecture~\ref{c2}(iii) on the absence of a variance scale outside the
coincidence matrices and the renormalized disjoint kernel,
Conjecture~\ref{c3}(ii) on the absence of a fourth covariance source,
Conjecture~\ref{c5}(ii) on the converse classification of rigid
coefficient sequences, Conjecture~\ref{c7}(iv) on the absence of mixed
additive-multiplicative terms, Conjecture~\ref{c9}(iii) on the rank-one
form of the exceptional-zero deformation, Conjecture~\ref{c11}(iii) on
the absence of a constraint beyond the relation subtorus,
Conjecture~\ref{c20}(ii) on the necessity of a geometric correspondence,
Conjecture~\ref{c22}(ii) on the absence of a further small-large
coupling, Conjecture~\ref{c23}(iii) on the absence of further colour
coupling, and Conjecture~\ref{c24}(ii) on the absence of a correlation
surviving the saddle-point tilt.
Refuting a converse layer identifies a missing invariant and leaves the
constructive layer of that conjecture standing, which is why the two are
never merged into one clause.

Two clauses are deductions rather than conjectures and are labelled as
such in place: Conjecture~\ref{c1}(iv), which derives Gaussianity and
Wick pairing from the connected cumulant hypothesis of
Conjecture~\ref{c1}(iii), and Conjecture~\ref{c17}(i), which derives the
subcritical total-variation bound from an effective Chebotarev estimate
under a Riemann hypothesis for the preimage fields.  Four clauses are
elementary and are proved by the argument stated with them:
Conjecture~\ref{c11}(i), Conjecture~\ref{c15}(i), Conjecture~\ref{c25}(iii)
and the identification of $q_p$ with a connecting map recorded in
Section~\ref{sec:notation}.  The constants $A(t)$ of
Conjecture~\ref{c6}, $\theta(H)$ of Conjecture~\ref{c4},
$\theta_{\mathrm{occ}}$ of Conjecture~\ref{c8}, $\kappa(\Gamma)$ of
Conjecture~\ref{c13}, $\delta_p$ of Conjecture~\ref{c19} and the mark
intensity of Conjecture~\ref{c21} are all derived from local, overlap,
Galois or sieve structure, and where a closed evaluation is not yet
available, as for $A(t)$, the derivation route is registered as the
stated programme of that conjecture.

Finally, the computational range available to the evidence notes is small
against every asymptotic asserted above.  The prime-pattern windows
overlap and carry no error bars, the occupancy moduli span less than one
decade of $\log q$, the finite-logarithm tests use fixed frequencies
only, the arboreal tests reach level four over ten effective maps, and
the factorization tests use samples of a few thousand values.  Each
evidence note says which clause it touches and which it leaves untested,
and none of them is offered as support for an asymptotic claim.

\end{document}
